\chapter[Structural DSGE Dynamics]{Structural Deterministic Dynamics in DSGE Filtering}
\label{ch:bf-structural-deterministic-dynamics}

\begin{mdframed}
\textbf{Production warning.}
Do not implement DSGE nonlinear filtering by silently treating every declared
state coordinate as an exchangeable noisy Gaussian coordinate.  Many DSGE
models contain a structural split between shock-driven exogenous states and
endogenous or accounting states that are deterministic conditional on lagged
states, controls, and current shocks.  This chapter proves the law-level reason:
for mixed structural models, the nonlinear prediction target is the pushforward
of the lagged filtering law and the current innovation law through the
structural transition.  A full-state additive-noise approximation is therefore
valid only when it preserves that law or is explicitly labeled and tested as an
approximation.
\end{mdframed}

\section{The Structural Split}
\label{sec:bf-dsge-structural-split}

The state vector in a DSGE likelihood is not merely a numerical vector.  It
usually carries economic timing.  In the language of the structural state
partition from Chapter~\ref{ch:bf-state-space-contracts}, a common DSGE
representation separates a declared stochastic block from a deterministic-
completion block.

\begin{definition}[Structural transition notation]
\label{def:bf-structural-transition-notation}
Throughout this chapter, write the full structural state as
$x_t=(m_t,k_t)$, where $m_t$ is the declared stochastic or exogenous block and
$k_t$ is the deterministic-completion block.  A mixed structural transition is
\begin{align}
  m_t &= T_m(m_{t-1}, \varepsilon_t; \theta), \label{eq:bf-m-transition}\\
  k_t &= T_k(k_{t-1}, m_{t-1}, m_t; \theta), \label{eq:bf-k-transition}\\
  x_t &= (m_t, k_t).
\end{align}
Equivalently, define the structural map
\begin{equation}
  F_\theta(x_{t-1},\varepsilon_t)
  =
  \left(
    T_m(m_{t-1},\varepsilon_t;\theta),
    T_k(k_{t-1},m_{t-1},
      T_m(m_{t-1},\varepsilon_t;\theta);\theta)
  \right).
  \label{eq:bf-structural-map-definition}
\end{equation}
Let $\pi_{t-1|t-1}^x$ denote the filtering law of $x_{t-1}$ and let
$\nu_\theta^\varepsilon$ denote the law of $\varepsilon_t$.  The pre-transition
law is
\[
  \pi_t^a = \pi_{t-1|t-1}^x \otimes \nu_\theta^\varepsilon
\]
when the current innovation is conditionally independent of the lagged
filtering law.  The predictive law of $x_t$ is
\[
  \pi_{t|t-1}^x(B)
  =
  \Pr_\theta(x_t\in B\mid y_{1:t-1})
\]
for measurable state sets $B$.  The pushforward of $\pi_t^a$ by $F_\theta$ is
the measure
\[
  (F_\theta)_\#\pi_t^a(B)
  =
  \pi_t^a\!\left(F_\theta^{-1}(B)\right).
\]
To place sigma points on a random variable means to choose weighted points
$\{a_t^{(j)},w_j^{(m)},w_j^{(c)}\}$ whose weighted mean and covariance match
the Gaussian approximation to that variable, then propagate
$x_t^{(j)}=F_\theta(a_t^{(j)})$ point by point.
\end{definition}

Thus, for mixed structural models,
\[
  \pi_{t|t-1}^x = (F_\theta)_\#\pi_t^a,
\]
or, in density notation when the densities exist,
\[
  p_\theta(x_t\mid y_{1:t-1})
  =
  \iint \delta\!\left(x_t - F_\theta(x_{t-1},\varepsilon_t)\right)
  p(\varepsilon_t)\,p_\theta(x_{t-1}\mid y_{1:t-1})
  \, d\varepsilon_t \, dx_{t-1}.
\]
In a perturbation solution, the same economics may also be represented by a
compact law of motion
\[
  x_t = h_x x_{t-1} + \eta \varepsilon_t
\]
at first order, with zero or rank-deficient innovation rows for the
deterministic-completion block.  The compact representation is useful, but it
must not erase the timing contract or the structural metadata described in
Chapter~\ref{ch:bf-api-design}.
In nonlinear filtering, the correct random input is therefore the innovation or
declared stochastic block.  The endogenous block is then generated by the
model-implied deterministic-completion map.  Here ``deterministic'' means no
independent innovation in that coordinate once $(x_{t-1},\varepsilon_t)$ is
fixed; it does not mean zero predictive variance conditional only on the lagged
filtering state.

This is the distinction emphasized in the DSGE filtering literature.  It is
also an instance of the generic BayesFilter structural-transition contract from
Chapter~\ref{ch:bf-state-space-contracts}: the backend integrates over the
declared stochastic block and completes the deterministic block pointwise.
Standard Bayesian DSGE expositions write the likelihood in state-space form
after the model solution has supplied a transition law for predetermined
variables and an observation equation for measured macroeconomic series
\citep{herbst2015bayesian,an2007bayesian}.  Pruned perturbation methods make
the same point in a different language: the first-order component is propagated
stochastically, while the higher-order correction is propagated by a separate
deterministic recursion driven by the lower-order state
\citep{kim2008calculating,andreasen2018pruned}.  The filtering backend must
honor the state-space law supplied by the structural solution, not merely the
dimension of the state vector.

\subsection{Literature Map for This Chapter}
\label{sec:bf-structural-literature-map}

The chapter uses four strands of literature, and each supports a different
claim.
\begin{enumerate}[label=(\roman*)]
  \item \textbf{State-space likelihoods and DSGE solutions.}
    Standard state-space and Bayesian DSGE references provide the filtering
    likelihood setup and the distinction between a structural model solution
    and the numerical recursion used to evaluate it
    \citep{durbin2012time,herbst2015bayesian,an2007bayesian}.
  \item \textbf{Pruned nonlinear DSGE state spaces.}
    Pruning papers show that nonlinear DSGE approximations naturally separate
    lower-order stochastic propagation from higher-order deterministic
    correction recursions
    \citep{kim2008calculating,andreasen2018pruned}.  That literature motivates
    the structural partition used here, but it does not by itself choose a
    sigma-point implementation.
  \item \textbf{Unscented and sigma-point filtering.}
    Julier and Uhlmann justify deterministic sigma points as a moment-matching
    approximation for nonlinear transformations, while van der Merwe's
    sigma-point Kalman-filter treatment supplies the dynamic state-space
    notation and augmented-noise filtering pattern used below
    \citep{julier1996general,julier1997new,vandermerwe2004sigma}.
  \item \textbf{Particle and factor backends.}
    Particle-filter references support the same law-level distinction in a
    simulation backend: sample the declared transition and weight by the
    observation density
    \citep{gordon1993novel,doucet2001sequential,andrieu2010particle}.
    Square-root and SVD chapters in this monograph address numerical
    factorization and derivative risk; they do not replace the structural-law
    requirement.
\end{enumerate}
The contribution of this chapter is to connect these strands: in a mixed DSGE
transition, sigma-point, SVD sigma-point, cubature, and particle methods should
approximate the pushforward law generated by the structural map
$F_\theta(x_{t-1},\varepsilon_t)$.

\section{Why Linear Kalman Filtering Can Hide the Problem}
\label{sec:bf-linear-kalman-hides-structural-split}

For a linear Gaussian state-space model,
\begin{align}
  x_t &= c + F x_{t-1} + u_t,\qquad u_t \sim \calN(0,Q),\\
  y_t &= a + H x_t + e_t,\qquad e_t \sim \calN(0,R),
\end{align}
the Kalman filter only needs the conditional mean and covariance of
$x_t\mid x_{t-1}$.  If a DSGE solution implies a singular or nearly singular
$Q$, the exact linear Gaussian likelihood remains well defined when the
collapsed representation preserves the same conditional law induced by the
structural transition.  Numerically stable implementations may then handle that
degenerate conditional Gaussian law through diffuse, square-root, solve-form,
or carefully regularized recursions as appropriate \citep{durbin2012time}.

This is why the following collapsed first-order representation can be
acceptable in an exact Kalman path:
\[
  Q = \eta \eta^\top,\qquad
  P_{t|t-1} = h_x P_{t-1|t-1} h_x^\top + Q.
\]
The key point is not that structural determinism disappears.  The key point is
that the collapsed linear representation still encodes the same one-step
conditional Gaussian law.  Deterministic-completion coordinates are not wrong
merely because they are included in $x_t$.  They obtain randomness through
their dependence on lagged stochastic coordinates, and their direct innovation
variance may be zero.  The exact linear Kalman recursion uses only the correct
first two conditional moments.

The danger is that this success can train the implementation culture into a
bad habit: it makes the state vector look like a homogeneous Gaussian vector.
That habit becomes wrong when the backend changes from exact linear filtering
to nonlinear sigma-point or particle filtering.  The relevant design decision
is then the integration space recorded in the filter metadata, not merely the
numerical dimension of $x_t$.

A related misunderstanding should be removed explicitly.  A deterministic-
completion coordinate need not have zero one-step predictive variance given the
lagged state.  It can inherit randomness through the current stochastic block.
In the Chapter~\ref{ch:bf-state-space-contracts} AR(2) lag-stack example, the
lag coordinate has no new innovation of its own but still remains random under
the lagged filtering distribution.  Likewise, in DSGE models with
\eqref{eq:bf-m-transition}--\eqref{eq:bf-k-transition}, $k_t$ may satisfy
\[
  \Var(k_t\mid x_{t-1}) > 0
\]
even though $k_t$ has no independent innovation once $(x_{t-1},\varepsilon_t)$
is fixed.

\section{Why Nonlinear Filters Must Respect the Structural Map}
\label{sec:bf-nonlinear-filters-need-structural-map}

Sigma-point and particle filters approximate a transition law, not just a
matrix covariance update.  In a generic nonlinear state-space model,
\[
  x_t = f_\theta(x_{t-1}, \varepsilon_t),\qquad
  y_t = g_\theta(x_t, e_t),
\]
the filter has to decide which variables are integrated over and which are
deterministic transforms.  In BayesFilter language, this is the choice of
integration space recorded in the filter metadata.  For DSGE models with
\eqref{eq:bf-m-transition}--\eqref{eq:bf-k-transition}, the natural nonlinear
prediction step is:
\begin{enumerate}[label=(\roman*)]
  \item represent the filtering distribution of the lagged structural state;
  \item draw, integrate, or choose quadrature points for the exogenous
    innovation $\varepsilon_t$ or current exogenous state $m_t$;
  \item compute each corresponding endogenous state $k_t$ using
    $T_k(k_{t-1},m_{t-1},m_t;\theta)$;
  \item evaluate the observation map on the resulting structurally valid
    points.
\end{enumerate}
The endogenous coordinates should not receive arbitrary independent
``sigma-point noise'' merely because they are entries of $x_t$.  If a sigma
point moves an accounting identity or predetermined endogenous state outside
the model-implied support of the structural transition, the filter is no longer
approximating the intended DSGE state-space model.  It is approximating a
different artificial model relative to the declared structural law.  This is
exactly the approximation boundary from Chapter~\ref{ch:bf-sigma-point-filters}:
quadrature may approximate the target law, but injecting additional stochastic
coordinates changes the target unless that change is declared and justified.

Particle-filter literature makes the same distinction operationally.  A
particle filter propagates particles through the state transition and weights
them by the observation density \citep{gordon1993novel,doucet2001sequential,
andrieu2010particle}.  Written in the structural notation, a basic propagation
and weighting step is
\[
  x_t^{(i)} = F_\theta\!\left(x_{t-1}^{(a_i)},\varepsilon_t^{(i)}\right),
  \qquad
  w_t^{(i)} \propto p_\theta\!\left(y_t\mid x_t^{(i)}\right).
\]
When part of the transition is deterministic, particles should be propagated by
sampling the declared stochastic block and then deterministically completing the
remaining coordinates.  Adding artificial noise to deterministic coordinates can
be used only as a declared proposal or regularization device, with the
proposal/target correction and approximation status stated explicitly.  It is
not the default filtering law.

\section{Reusable BayesFilter Structural Filtering Step}
\label{sec:bf-structural-filter-step-algorithm}

The doctrine above can be stated as a reusable backend contract.  At every time
step, the filter should treat the structural transition as the object being
approximated:
\[
  x_t = F_\theta(x_{t-1},\varepsilon_t),
  \qquad
  y_t \sim p_\theta(y_t\mid x_t).
\]
The filtered state is still the full state $x_t$.  The integration variable is
the pre-transition uncertainty that generates the predictive law.

\begin{mdframed}
\textbf{Algorithm: \texttt{structural\_filter\_step}.}

\begin{enumerate}[label=(\arabic*)]
  \item Start from the filtered approximation to
    $p_\theta(x_{t-1}\mid y_{1:t-1})$ and the declared innovation law
    $p_\theta(\varepsilon_t)$.
  \item Form sigma points, quadrature points, particles, or linearization
    points in the declared integration space:
    $(x_{t-1},\varepsilon_t)$, $(s_{t-1},d_{t-1},\varepsilon_t)$, or an
    equivalent stochastic-state parameterization.
  \item For each point, call the structural transition map:
    \[
      s_t^{(j)}=T_s(s_{t-1}^{(j)},\varepsilon_t^{(j)};\theta),
      \qquad
      d_t^{(j)}=T_d(d_{t-1}^{(j)},s_{t-1}^{(j)},s_t^{(j)};\theta).
    \]
    Deterministic-completion coordinates are outputs of this map, not
    independent new stochastic coordinates.
  \item Pack $x_t^{(j)}=\operatorname{pack}(s_t^{(j)},d_t^{(j)},a_t^{(j)})$
    and evaluate the observation map or density on these structurally valid
    points.
  \item Apply the selected update: exact Kalman, Gaussian sigma-point closure,
    particle weighting/resampling, or another declared backend.
  \item Return the likelihood contribution, filtered approximation, and
    metadata recording integration space, deterministic-completion policy,
    approximation label, derivative policy, and finite-failure diagnostics.
\end{enumerate}
\end{mdframed}

This algorithm is intentionally backend-neutral.  For a linear Gaussian Kalman
backend, the structural transition may collapse to matrices $(F,Q)$ as long as
the same conditional Gaussian law is preserved.  For UKF, CKF, SVD sigma-point,
and particle backends, the pointwise structural transition is the default.  A
mixed-model full-state Gaussian approximation may still be offered as a
diagnostic or legacy path, but it must carry an approximation label and must be
tested against structural references.

\section[Degenerate Transitions]{Degenerate Transitions Are Structural}
\label{sec:bf-degenerate-transitions-structural}

Degenerate transition covariance is not an exceptional corner case in DSGE
models.  It is often the mathematical expression of timing, identities, and
endogenous accumulation.  Capital, debt, lagged output growth, risk-premium
terms, and other predetermined variables may be deterministic conditional on
the previous state and current shocks.  A full covariance matrix with small positive diagonal padding can keep linear
algebra routines alive, but it does not by itself define the intended economic
transition.  BayesFilter therefore separates three ideas that are often
confused in practice:
\begin{enumerate}[label=(\roman*)]
  \item an exact degenerate transition law that already matches the structural
    model;
  \item numerical regularization used to evaluate that law stably; and
  \item an altered transition law that injects new stochastic degrees of
    freedom into deterministic-completion coordinates.
\end{enumerate}
Only the first two preserve the original structural model automatically.

The implementation policy is therefore:
\begin{itemize}
  \item distinguish structural determinism from numerical regularization;
  \item never let a nugget or PSD projection change the declared stochastic
    dimension without a written approximation label;
  \item expose deterministic blocks in diagnostics, rather than burying them
    under a full-rank covariance;
  \item test degenerate and rank-deficient transition cases directly.
\end{itemize}

This point is separate from the SVD spectral-gradient issue in
Chapter~\ref{ch:bf-svd-sigma-point}.  SVD factors may improve value-side
robustness for nearly singular covariances, and an SVD sigma-point filter can
be the right numerical backend for a structurally declared sigma-point rule.
But SVD is a factorization choice, not a state-law declaration.  It repairs the
reviewer's degeneracy concern only if the sigma-point backend factors the
covariance of the correct integration variable, such as
$(x_{t-1},\varepsilon_t)$, and then completes deterministic coordinates
pointwise.  If the backend first invents a full-rank post-transition covariance
for deterministic-completion coordinates, an SVD factor merely represents that
altered law more stably.

\section{Standard UKF, Structural UKF, and the Predictive Law}
\label{sec:bf-ukf-structural-overview}

\subsection{The Original Unscented-Transform Pattern}
\label{sec:bf-original-ut-pattern}

Julier and Uhlmann's original unscented-transform construction starts from a
random variable $X\in\mathbb{R}^L$ with mean $\bar x$ and covariance $P_x$ and
approximates the moments of $Y=g(X)$ by propagating deterministic sigma points
through $g$ \citep{julier1996general,julier1997new}.  The
$(\alpha,\beta,\kappa)$ scaled rule used for the formulas below belongs to the
later sigma-point Kalman filter convention \citep{vandermerwe2004sigma}.  In
the notation used here, one convenient scaled UKF rule forms
\begin{align}
  \chi_0 &= \bar x,\\
  \chi_i &= \bar x+\left[\sqrt{(L+\lambda)P_x}\right]_i,
    \qquad i=1,\ldots,L,\\
  \chi_{i+L} &= \bar x-\left[\sqrt{(L+\lambda)P_x}\right]_i,
    \qquad i=1,\ldots,L,
    \label{eq:bf-scaled-ukf-points}
\end{align}
where $[\sqrt{(L+\lambda)P_x}]_i$ denotes the $i$th column or row of the chosen
matrix square root.  With
\[
  \lambda=\alpha^2(L+\kappa)-L,
\]
the usual scaled weights are
\begin{align}
  w_0^{(m)} &= \frac{\lambda}{L+\lambda},&
  w_0^{(c)} &= \frac{\lambda}{L+\lambda}+1-\alpha^2+\beta,\\
  w_i^{(m)} &= w_i^{(c)} = \frac{1}{2(L+\lambda)},
  & i&=1,\ldots,2L.
  \label{eq:bf-scaled-ukf-weights}
\end{align}
The propagated points are $\upsilon_i=g(\chi_i)$, and the transformed moments
are approximated by
\begin{align}
  \bar y &\approx \sum_{i=0}^{2L} w_i^{(m)}\upsilon_i,\\
  P_y &\approx
    \sum_{i=0}^{2L} w_i^{(c)}
      (\upsilon_i-\bar y)(\upsilon_i-\bar y)^\top .
  \label{eq:bf-ut-transformed-moments}
\end{align}
The early unscented-transform report analyzes this construction by Taylor
expansion: the method matches the input mean and covariance and pushes the
resulting deterministic ensemble through the nonlinear map, yielding lower-order
moment accuracy that improves on first-order linearization under the smoothness
conditions stated there \citep{julier1996general}.  Van der Merwe's
sigma-point Kalman filter treatment uses the same family logic: sigma points
are selected to capture the important first and second moments of the prior
Gaussian random variable and then weighted after nonlinear propagation
\citep{vandermerwe2004sigma}.

This chapter uses that scaled UKF convention for exposition.  Later UKF, CDKF,
and square-root variants differ in point placement, weights, and factorization
details, but the structural question is invariant: which random variable is the
sigma-point rule approximating?  If the wrong variable is chosen, the transform
may approximate the wrong law accurately.

In the additive-noise filtering setting, the unscented-transform idea is usually
applied to an augmented Gaussian variable containing the prior state and current
additive-noise variables.  The sigma points therefore approximate the variables
that generate the next predicted state and predicted observation.  For mixed
structural DSGE transitions, that same principle is exactly why the sigma-point
variable is $(x_{t-1},\varepsilon_t)$ or an equivalent pre-transition
parameterization, not an artificially perturbed post-transition full state.

\subsection{Standard Additive-Noise UKF Algorithm}
\label{sec:bf-standard-ukf-algorithm}

To avoid confusion, write the additive-noise UKF pattern explicitly.  Suppose
the model is declared as
\[
  x_t = f_\theta(x_{t-1},u_t),
  \qquad
  y_t = h_\theta(x_t,e_t),
\]
with independent Gaussian state innovation $u_t\sim\calN(0,Q_t)$ and
measurement innovation $e_t\sim\calN(0,R_t)$.  Given a previous Gaussian
filtering approximation, one additive-noise UKF step is:
\begin{enumerate}[label=\textbf{U\arabic*.}]
  \item \textbf{Inputs.}  Start from
    \[
      x_{t-1}\mid y_{1:t-1} \approx \calN(m_{t-1|t-1},P_{t-1|t-1}).
    \]
    The implementation receives $f_\theta$, $h_\theta$, $Q_t$, $R_t$, the
    observed value $y_t$, sigma-point parameters
    $(\alpha,\beta,\kappa)$, and a declared covariance-factorization policy.
  \item \textbf{Augmented variable.}  Form
    \[
      b_t=(x_{t-1},u_t,e_t),
      \qquad
      b_t\mid y_{1:t-1}\approx \calN(\mu_b,P_b),
    \]
    with
    \[
      \mu_b=(m_{t-1|t-1},0,0),\qquad
      P_b=\operatorname{blockdiag}(P_{t-1|t-1},Q_t,R_t).
    \]
    If a variant handles measurement noise by adding $R_t$ after observation
    propagation, omit $e_t$ from $b_t$ and add $R_t$ in step U6.
  \item \textbf{Sigma points.}  Generate
    $\{b_t^{(j)},w_j^{(m)},w_j^{(c)}\}_{j=0}^{2L_b}$ from
    $(\mu_b,P_b)$ using \eqref{eq:bf-scaled-ukf-points}--%
    \eqref{eq:bf-scaled-ukf-weights}.  Write
    $b_t^{(j)}=(x_{t-1}^{(j)},u_t^{(j)},e_t^{(j)})$.
  \item \textbf{State propagation.}  Propagate
    \[
      x_t^{(j)}=f_\theta(x_{t-1}^{(j)},u_t^{(j)}).
    \]
    Compute
    \[
      \hat x_{t|t-1}=\sum_j w_j^{(m)}x_t^{(j)},\qquad
      P_{xx,t}=\sum_j w_j^{(c)}
      (x_t^{(j)}-\hat x_{t|t-1})(x_t^{(j)}-\hat x_{t|t-1})^\top .
    \]
  \item \textbf{Observation propagation.}  Form
    \[
      z_t^{(j)}=h_\theta(x_t^{(j)},e_t^{(j)})
    \]
    or $z_t^{(j)}=h_\theta(x_t^{(j)},0)$ with $R_t$ added below, depending on
    the chosen variant.
  \item \textbf{Innovation moments.}  Compute
    \begin{align}
      \hat z_t &= \sum_j w_j^{(m)}z_t^{(j)},\\
      S_t &=
        \sum_j w_j^{(c)}(z_t^{(j)}-\hat z_t)(z_t^{(j)}-\hat z_t)^\top
        + R_t^{\mathrm{add}},\\
      P_{xz,t} &=
        \sum_j w_j^{(c)}
        (x_t^{(j)}-\hat x_{t|t-1})(z_t^{(j)}-\hat z_t)^\top,
    \end{align}
    where $R_t^{\mathrm{add}}=0$ if measurement noise was included in $b_t$ and
    $R_t^{\mathrm{add}}=R_t$ otherwise.
  \item \textbf{Gaussian update and likelihood.}  With
    $v_t=y_t-\hat z_t$,
    \begin{align}
      K_t &= P_{xz,t}S_t^{-1},\\
      \hat x_{t|t} &= \hat x_{t|t-1}+K_tv_t,\\
      P_{t|t} &= P_{xx,t}-K_tS_tK_t^\top,\\
      \ell_t &=
      -\frac{1}{2}
      \left[
        d_y\log(2\pi)+\log|S_t|+v_t^\top S_t^{-1}v_t
      \right].
    \end{align}
  \item \textbf{Return metadata.}  Return
    $(\ell_t,\hat x_{t|t},P_{t|t})$ with metadata recording the augmented
    variable $b_t$, the factorization method, sigma-point parameters, whether
    measurement noise was augmented or added, and any PSD/jitter/finiteness
    interventions.
\end{enumerate}
This is the textbook pattern many readers have in mind.  It is perfectly
natural when the declared state transition itself is an additive-noise model on
the full state.

\subsection{Structural UKF Algorithm}
\label{sec:bf-structural-ukf-algorithm}

For a mixed structural model, the update algebra remains the same, but the
sigma-point variable is chosen differently because the predictive law is
generated differently.  For the structural transition
\eqref{eq:bf-m-transition}--\eqref{eq:bf-k-transition}, a structural UKF step
is:
\begin{enumerate}[label=\textbf{S\arabic*.}]
  \item \textbf{Inputs.}  Start from the filtering approximation
    \[
      x_{t-1}\mid y_{1:t-1} \approx \calN(\mu_{t-1},P_{t-1}).
    \]
    The implementation receives $T_m$, $T_k$, the observation map
    $h_\theta$, the innovation law
    $\varepsilon_t\sim\calN(0,\Sigma_{\varepsilon,t})$, observation covariance
    $R_t$, the observed value $y_t$, sigma-point parameters, and a
    factorization policy.
  \item \textbf{Pre-transition variable.}  Form
    \[
      a_t = (x_{t-1},\varepsilon_t),
      \qquad
      a_t\mid y_{1:t-1}\approx\calN(\mu_a,P_a),
    \]
    with
    \[
      \mu_a=(\mu_{t-1},0),\qquad
      P_a=\operatorname{blockdiag}(P_{t-1},\Sigma_{\varepsilon,t}).
    \]
    This is the sigma-point variable because
    $\pi_{t|t-1}^x=(F_\theta)_\#\pi_t^a$.
  \item \textbf{Sigma points.}  Generate
    $\{a_t^{(j)},w_j^{(m)},w_j^{(c)}\}_{j=0}^{2L_a}$ from
    $(\mu_a,P_a)$.  Write
    $a_t^{(j)}=(m_{t-1}^{(j)},k_{t-1}^{(j)},\varepsilon_t^{(j)})$.
  \item \textbf{Stochastic-block propagation.}  For each point compute
    \[
      m_t^{(j)}
      =
      T_m(m_{t-1}^{(j)},\varepsilon_t^{(j)};\theta).
    \]
  \item \textbf{Deterministic completion.}  For the same point compute
    \[
      k_t^{(j)}
      =
      T_k(k_{t-1}^{(j)},m_{t-1}^{(j)},m_t^{(j)};\theta),
      \qquad
      x_t^{(j)}=(m_t^{(j)},k_t^{(j)}).
    \]
    There is no independent sigma-point coordinate for $k_t$ unless the model
    declares one.
  \item \textbf{State moments.}  Compute
    \[
      \hat x_{t|t-1}=\sum_j w_j^{(m)}x_t^{(j)},\qquad
      P_{xx,t}=\sum_j w_j^{(c)}
      (x_t^{(j)}-\hat x_{t|t-1})(x_t^{(j)}-\hat x_{t|t-1})^\top .
    \]
  \item \textbf{Observation points and moments.}  Evaluate
    $z_t^{(j)}=h_\theta(x_t^{(j)},0)$ or the appropriate noise-augmented
    observation variant, then compute
    \begin{align}
      \hat z_t &= \sum_j w_j^{(m)}z_t^{(j)},\\
      S_t &=
        \sum_j w_j^{(c)}(z_t^{(j)}-\hat z_t)(z_t^{(j)}-\hat z_t)^\top + R_t,\\
      P_{xz,t} &=
        \sum_j w_j^{(c)}
        (x_t^{(j)}-\hat x_{t|t-1})(z_t^{(j)}-\hat z_t)^\top .
    \end{align}
  \item \textbf{Gaussian update and likelihood.}  Apply exactly the update from
    U7:
    \[
      v_t=y_t-\hat z_t,\qquad
      K_t=P_{xz,t}S_t^{-1},\qquad
      \hat x_{t|t}=\hat x_{t|t-1}+K_tv_t,
    \]
    with $P_{t|t}=P_{xx,t}-K_tS_tK_t^\top$ and the same Gaussian innovation
    log-likelihood formula.
  \item \textbf{Return metadata.}  Return the likelihood contribution and
    filtered Gaussian approximation with metadata recording
    \code{integration\_space} equal to $(x_{t-1},\varepsilon_t)$,
    \code{deterministic\_completion} equal to \code{pointwise}, structural
    identity diagnostics, sigma-point parameters, factorization method, and any
    approximation label.
\end{enumerate}
Structural UKF is not a different Kalman-gain formula.  It is the same Gaussian
update algebra applied to a predictive sigma-point cloud that respects the
structural transition law.

\subsection{UKF in Sequential-Monte-Carlo Kernel Notation}
\label{sec:bf-ukf-kernel-notation}

The same point can be expressed in the probability-kernel notation used in
\citet{chopin2020introduction}.  A state-space model is specified by an initial
law $P_{\theta,0}(dx_0)$, transition kernels
$P_{\theta,t}(x_{t-1},dx_t)$, and observation densities
$f_{\theta,t}(y_t\mid x_t)$.  In the bootstrap Feynman--Kac formalism, the
transition kernel is
\[
  M_{\theta,t}(x_{t-1},dx_t)=P_{\theta,t}(x_{t-1},dx_t),
\]
and the potential is
\[
  G_{\theta,t}(x_t)=f_{\theta,t}(y_t\mid x_t).
\]
If $\eta_{t-1}$ denotes the filtering law of $x_{t-1}$ given
$y_{1:t-1}$, the exact one-step recursion is
\begin{align}
  \xi_t(dx_t)
    &=
    \int \eta_{t-1}(dx_{t-1})
      P_{\theta,t}(x_{t-1},dx_t),
    \label{eq:bf-kernel-exact-prediction}\\
  \lambda_t
    &=
    \xi_t(G_{\theta,t})
    =
    \int G_{\theta,t}(x_t)\,\xi_t(dx_t),
    \label{eq:bf-kernel-exact-likelihood-factor}\\
  \eta_t(dx_t)
    &=
    \frac{G_{\theta,t}(x_t)\xi_t(dx_t)}{\xi_t(G_{\theta,t})}.
    \label{eq:bf-kernel-exact-update}
\end{align}
Equations~\eqref{eq:bf-kernel-exact-prediction}--%
\eqref{eq:bf-kernel-exact-update} are the filtering target.  A particle filter
approximates the measures in these equations by random weighted particles.  A
UKF instead approximates the same prediction/update recursion by deterministic
quadrature and Gaussian projection.

For the mixed structural transition in
\eqref{eq:bf-structural-map-definition}, the transition kernel is the pushforward
kernel
\begin{equation}
  P_{\theta,t}^{\mathrm{str}}(x,B)
  =
  \int
  \mathbf 1\!\left\{F_\theta(x,\varepsilon)\in B\right\}
  \nu_{\theta,t}^{\varepsilon}(d\varepsilon)
  =
  \int
  \delta_{F_\theta(x,\varepsilon)}(B)
  \nu_{\theta,t}^{\varepsilon}(d\varepsilon),
  \label{eq:bf-structural-transition-kernel}
\end{equation}
for measurable state sets $B$.  This is a perfectly valid probability kernel
even when it is singular with respect to Lebesgue measure.  Kernel notation is
therefore useful here: it does not force a degenerate structural transition to
pretend that it has a full-dimensional transition density.  The predictive
integral in \eqref{eq:bf-kernel-exact-prediction} becomes, for any test function
$\psi$,
\begin{equation}
  \xi_t^{\mathrm{str}}(\psi)
  =
  \iint
  \psi\!\left(F_\theta(x_{t-1},\varepsilon_t)\right)
  \eta_{t-1}(dx_{t-1})
  \nu_{\theta,t}^{\varepsilon}(d\varepsilon_t).
  \label{eq:bf-structural-kernel-prediction-test-function}
\end{equation}
Thus the kernel version of the structural statement is:
\[
  P_{\theta,t}=P_{\theta,t}^{\mathrm{str}},
  \qquad
  M_{\theta,t}=P_{\theta,t}^{\mathrm{str}},
  \qquad
  G_{\theta,t}(x)=f_{\theta,t}(y_t\mid x).
\]

Now suppose the incoming filter is approximated by a Gaussian
$\widehat\eta_{t-1}=\calN(\widehat m_{t-1},\widehat P_{t-1})$.  Let
\[
  \widehat\alpha_t(da)
  =
  \widehat\eta_{t-1}(dx_{t-1})
  \nu_{\theta,t}^{\varepsilon}(d\varepsilon_t),
  \qquad
  a=(x_{t-1},\varepsilon_t).
\]
The structural UKF replaces integrals against $\widehat\alpha_t$ by the
deterministic sigma-point rule
\[
  \mathcal U_t^{(r)}[\psi]
  =
  \sum_j w_j^{(r)}\psi(a_t^{(j)}),
  \qquad r\in\{m,c\}.
\]
With
\[
  x_t^{(j)}=F_\theta(a_t^{(j)}),
  \qquad
  z_t^{(j)}=h_\theta(x_t^{(j)}),
\]
the UKF predictive and observation moments are
\begin{align}
  \widehat m_t^-
    &=
    \sum_j w_j^{(m)}x_t^{(j)},\\
  \widehat P_t^-
    &=
    \sum_j w_j^{(c)}
    (x_t^{(j)}-\widehat m_t^-)(x_t^{(j)}-\widehat m_t^-)^\top,
    \label{eq:bf-kernel-ukf-state-moments}\\
  \widehat z_t
    &=
    \sum_j w_j^{(m)}z_t^{(j)},\\
  \widehat S_t
    &=
    \sum_j w_j^{(c)}
    (z_t^{(j)}-\widehat z_t)(z_t^{(j)}-\widehat z_t)^\top
    +R_t,\\
  \widehat P_{xz,t}
    &=
    \sum_j w_j^{(c)}
    (x_t^{(j)}-\widehat m_t^-)(z_t^{(j)}-\widehat z_t)^\top .
    \label{eq:bf-kernel-ukf-observation-moments}
\end{align}
Equivalently, UKF builds the Gaussian joint approximation
\begin{equation}
  \widehat J_t(dx_t,dy_t)
  =
  \calN\!\left(
    \begin{bmatrix}x_t\\ y_t\end{bmatrix};
    \begin{bmatrix}\widehat m_t^-\\ \widehat z_t\end{bmatrix},
    \begin{bmatrix}
      \widehat P_t^- & \widehat P_{xz,t}\\
      \widehat P_{xz,t}^\top & \widehat S_t
    \end{bmatrix}
  \right)dx_t\,dy_t .
  \label{eq:bf-kernel-ukf-joint-gaussian}
\end{equation}
The UKF likelihood factor and update are then the marginal and conditional
Gaussian objects induced by \eqref{eq:bf-kernel-ukf-joint-gaussian}:
\begin{align}
  \widehat\lambda_t^{\mathrm{UKF}}
    &=
    \varphi_{d_y}(y_t;\widehat z_t,\widehat S_t),
    \label{eq:bf-kernel-ukf-likelihood-factor}\\
  \widehat\eta_t^{\mathrm{UKF}}(dx_t)
    &=
    \widehat J_t(dx_t\mid y_t) \nonumber\\
    &=
    \calN\!\left(
      dx_t;\,
      \widehat m_t^-+
      \widehat P_{xz,t}\widehat S_t^{-1}(y_t-\widehat z_t),
      \widehat P_t^--
      \widehat P_{xz,t}\widehat S_t^{-1}\widehat P_{xz,t}^\top
    \right).
    \label{eq:bf-kernel-ukf-gaussian-update}
\end{align}
The corresponding contribution to the log likelihood is
$\log\widehat\lambda_t^{\mathrm{UKF}}$.
For actual stochastic-volatility rows, augmentation is admissible only when
marginalizing the augmented variable reproduces the same observation density
intended by the inference problem.  In the transformed actual-SV setting that
means preserving the exact law
\(
  z_t=\log(y_t^2),
  \ z_t-\log(\beta^2)-h_t\sim\log(\chi_1^2)
\).
The previous actual-SV Gaussian-closure framing failed this condition: it used
augmentation to construct a different innovation scalar rather than the intended
normalizer.  That construction may still define a Gaussian-projection filter, but
it does not define the correct actual-SV likelihood target.
If measurement noise is augmented rather than added, replace
$a=(x_{t-1},\varepsilon_t)$ by
$a=(x_{t-1},\varepsilon_t,e_t)$, set
$z_t^{(j)}=h_\theta(F_\theta(x_{t-1}^{(j)},\varepsilon_t^{(j)}),e_t^{(j)})$,
and omit the extra $R_t$ term in \eqref{eq:bf-kernel-ukf-observation-moments}.

This notation also states the equivalence condition for collapsed SVD
constructions.  If a full-state degenerate SVD implementation uses another
kernel
\[
  \widetilde P_{\theta,t}(x,B)
  =
  \int
  \delta_{\widetilde F_\theta(x,\eta)}(B)
  \widetilde\nu_{\theta,t}(d\eta),
\]
then it is a law-equivalent implementation of the structural model only when
\begin{equation}
  \widetilde P_{\theta,t}(x,B)
  =
  P_{\theta,t}^{\mathrm{str}}(x,B)
  \qquad
  \text{for the relevant }x\text{ and measurable }B,
  \label{eq:bf-kernel-equivalence-condition}
\end{equation}
or at least when the two kernels induce the same predictive measure from the
current filtering approximation.  Equality of a covariance matrix, or existence
of an SVD factor for a singular covariance, is not by itself equality of
probability kernels.  This is the kernel-notation version of the structural
warning in this chapter.

\subsection{Why These Algorithms Differ}
\label{sec:bf-ukf-algorithm-contrast}

\begin{center}
\begin{tabular}{p{0.20\linewidth}p{0.34\linewidth}p{0.34\linewidth}}
\toprule
\textbf{Object} & \textbf{Standard additive-noise UKF} & \textbf{Structural UKF}\\
\midrule
Filtered object & Full state $x_t$ & Full state $x_t$\\
Sigma-point variable & Declared augmented additive-noise variable
$(x_{t-1},u_t,e_t)$ & Pre-transition structural variable
$(x_{t-1},\varepsilon_t)$\\
Direct process noise & Whatever the additive model declares & Only the declared
stochastic block\\
Deterministic coordinates & None unless constraints are declared & Completed by
$T_k$ pointwise\\
Target law & Additive-noise predictive law & Structural pushforward law\\
Failure if misused & Moment approximation error & Model-law error plus moment
approximation error\\
SVD/square-root role & Covariance factorization & Covariance factorization, not
structural repair\\
\bottomrule
\end{tabular}
\end{center}

The update equations are shared.  The exact difference is upstream: which
random variable receives sigma points and which map generates the predictive
state cloud.

\subsection{Why the Sigma-Point Variable Uses Pre-Transition Uncertainty}
\label{sec:bf-structural-ukf-proposition}

\begin{proposition}[Predictive pushforward and sigma-point space]
\label{prop:bf-structural-ukf-pushforward}
Under the structural transition
\begin{align}
  m_t &= T_m(m_{t-1},\varepsilon_t),
    \label{eq:bf-structural-ukf-prop-m}\\
  k_t &= T_k(k_{t-1},m_{t-1},m_t),
    \label{eq:bf-structural-ukf-prop-k}\\
  x_t &= (m_t,k_t),
    \label{eq:bf-structural-ukf-prop-state}
\end{align}
with innovation law $p(\varepsilon_t)$ independent of
$x_{t-1}\mid y_{1:t-1}$, the following hold:
\begin{enumerate}[label=(\roman*)]
  \item the one-step predictive law $p(x_t\mid y_{1:t-1})$ is the pushforward
    of the joint law of $(x_{t-1},\varepsilon_t)$ under the structural map;
  \item a UKF that approximates this predictive law should place sigma points on
    $(x_{t-1},\varepsilon_t)$, or on an equivalent parameterization of the same
    pre-transition uncertainty;
  \item a naive full-state additive-noise UKF targets the same predictive law
    only if its induced predictive law coincides with that pushforward law; if
    it adds a direct perturbation to $k_t$, it targets a different law.
\end{enumerate}
\end{proposition}

\begin{proof}
Define the structural map
\begin{equation}
  F(x_{t-1},\varepsilon_t)
  =
  \bigl(T_m(m_{t-1},\varepsilon_t),
  T_k(k_{t-1},m_{t-1},T_m(m_{t-1},\varepsilon_t))\bigr).
  \label{eq:bf-structural-ukf-prop-map}
\end{equation}
Because the current innovation is independent of the lagged filtering law, the
joint pre-transition uncertainty factorizes as
\begin{equation}
  p(x_{t-1},\varepsilon_t\mid y_{1:t-1})
  =
  p(x_{t-1}\mid y_{1:t-1})p(\varepsilon_t).
  \label{eq:bf-structural-ukf-prop-factorization}
\end{equation}
Let $\varphi$ be any bounded test function on the state space of $x_t$.  Then
\begin{equation}
  \E\!\left[\varphi(x_t)\mid y_{1:t-1}\right]
  =
  \iint
  \varphi\!\left(F(x_{t-1},\varepsilon_t)\right)
  p(x_{t-1}\mid y_{1:t-1})p(\varepsilon_t)
  \, d\varepsilon_t \, dx_{t-1}.
  \label{eq:bf-structural-ukf-prop-pushforward}
\end{equation}
Equation~\eqref{eq:bf-structural-ukf-prop-pushforward} is exactly the
expectation of $\varphi$ under the pushforward of the joint law of
$(x_{t-1},\varepsilon_t)$ through $F$, which proves part~(i).

Moreover, under the structural transition,
\begin{equation}
  x_t = F(x_{t-1},\varepsilon_t)
  \qquad \text{almost surely.}
  \label{eq:bf-structural-ukf-prop-deterministic-completion}
\end{equation}
So once $(x_{t-1},\varepsilon_t)$ is fixed, both $m_t$ and $k_t$ are fixed.
There is no independent new perturbation attached to $k_t$ under the intended
model.  Therefore any sigma-point rule intended to approximate
$p(x_t\mid y_{1:t-1})$ should be applied to the pre-transition uncertainty
variables that generate that law, namely $(x_{t-1},\varepsilon_t)$ or an
exactly equivalent parameterization.  This proves part~(ii).

Finally, if a full-state additive-noise UKF omits $\varepsilon_t$, then it
misses genuine current-shock uncertainty.  If it instead adds a direct
perturbation to $k_t$, then it enlarges the stochastic space and defines a
different predictive law unless that perturbation is degenerate in exactly the
way required to reproduce the same pushforward distribution in
\eqref{eq:bf-structural-ukf-prop-pushforward}.  This proves part~(iii).
\end{proof}

The proposition does not say that the filter stops targeting $x_t$.  The
filtered state remains $x_t$.  The point is that the sigma-point variable is the
pre-transition uncertainty whose pushforward generates the predictive law of
$x_t$.  This is analogous to process-noise augmentation in a standard additive-
noise UKF, but in the mixed structural setting the distinction is sharper:
$\varepsilon_t$ is a genuine new source of uncertainty, while $k_t$ is a
structural completion of that uncertainty rather than an independently shocked
coordinate.

\begin{proposition}[Local structural UKF Taylor accuracy]
\label{prop:bf-structural-ukf-ut-accuracy}
Let $A_t=(x_{t-1},\varepsilon_t)$ have a Gaussian approximation with mean
$\mu_a$ and covariance $P_a$, and let
$X_t=F_\theta(A_t)$ where $F_\theta$ is the structural map in
\eqref{eq:bf-structural-map-definition}.  Write
$P_a=\eta^2\Sigma_a$, with $\Sigma_a$ fixed and $\eta$ measuring the local
scale of the input uncertainty.  Suppose the sigma-point deviations
$\xi_j=a_t^{(j)}-\mu_a$ satisfy
\[
  \sum_j w_j^{(m)}\xi_j=0,
  \qquad
  \sum_j w_j^{(m)}\xi_j\xi_j^\top=P_a,
  \qquad
  \sum_j w_j^{(c)}\xi_j\xi_j^\top=P_a,
\]
and suppose $F_\theta$ is three-times continuously differentiable on a
neighborhood containing the relevant input mass and sigma points.  Let $J$ be
the Jacobian of $F_\theta$ at $\mu_a$, and let $H_r$ be the Hessian of the
$r$th component of $F_\theta$ at $\mu_a$.  Then the structural UKF prediction
matches the transformed mean through the quadratic Taylor term:
\begin{align}
  \E[X_{t,r}\mid y_{1:t-1}]
    &=
    F_{\theta,r}(\mu_a)
    +\frac{1}{2}\operatorname{tr}(H_rP_a)
    +O(\eta^3),\\
  \hat x_{t|t-1,r}^{\mathrm{UKF}}
    &=
    F_{\theta,r}(\mu_a)
    +\frac{1}{2}\operatorname{tr}(H_rP_a)
    +O(\eta^3).
    \label{eq:bf-structural-ukf-second-order-mean}
\end{align}
It also matches the leading second-order covariance term:
\begin{align}
  \operatorname{Cov}(X_t\mid y_{1:t-1})
    &=
    JP_aJ^\top+O(\eta^3),\\
  P_{xx,t}^{\mathrm{UKF}}
    &=
    JP_aJ^\top+O(\eta^3).
    \label{eq:bf-structural-ukf-second-order-cov}
\end{align}
Thus, under the local small-uncertainty Taylor expansion used here, the
structural UKF matches the transformed mean through the quadratic term and
matches the leading covariance term for the Gaussian approximation on the
correct input law $(x_{t-1},\varepsilon_t)$.  This is a local expansion
statement for the selected structural input law, not a blanket exactness theorem
for nonlinear structural transitions.  It does not apply to a naive
post-transition full-state sigma-point cloud unless that cloud induces the same
law as $(F_\theta)_\#\pi_t^a$.
\end{proposition}

\begin{proof}
By Proposition~\ref{prop:bf-structural-ukf-pushforward}, the structural
predictive law of $x_t$ is the law of $F_\theta(A_t)$.  The structural UKF
therefore applies the unscented-transform rule to the transformation
$g=F_\theta$ and the input variable $A_t$.

For a component $r$, Taylor expansion around $\mu_a$ gives
\[
  F_{\theta,r}(\mu_a+\delta)
  =
  F_{\theta,r}(\mu_a)
  +J_r\delta
  +\frac{1}{2}\delta^\top H_r\delta
  +O(\|\delta\|^3).
\]
Taking expectation under the Gaussian approximation for
$A_t-\mu_a$ yields the first line of
\eqref{eq:bf-structural-ukf-second-order-mean}, because
$\E[A_t-\mu_a]=0$ and
$\E[(A_t-\mu_a)(A_t-\mu_a)^\top]=P_a$.  Taking the weighted sum over sigma
points gives the second line of
\eqref{eq:bf-structural-ukf-second-order-mean} by the same first- and
second-moment identities for $\xi_j$.

For covariance, subtract the corresponding mean expansion.  The leading
centered term is $J(A_t-\mu_a)$ for the true law and $J\xi_j$ for the sigma
points.  Their second moments are both $JP_aJ^\top$ by moment matching.  The
quadratic Taylor terms have size $O(\eta^2)$, so their covariance contribution
with the leading linear term is $O(\eta^3)$ and their self-covariance
contribution is $O(\eta^4)$.  This proves
\eqref{eq:bf-structural-ukf-second-order-cov}.  These are the local
unscented-transform moment calculations for the selected sigma-point rule
\citep{julier1996general,vandermerwe2004sigma}, applied to the structural
input variable.  If a naive full-state cloud is constructed from a different
input law, the same UT calculation applies only to that different law; it gives
no accuracy guarantee for the structural pushforward target.
\end{proof}

\section{Worked Example A: Nonlinear Structural Transition}
\label{sec:bf-nonlinear-structural-transition-example}

Let
\begin{align}
  m_t &= \rho m_{t-1}+\sigma\varepsilon_t,
    & \varepsilon_t &\sim \calN(0,1), \label{eq:bf-ukf-example-m}\\
  k_t &= \phi k_{t-1}+\gamma m_t^2, \label{eq:bf-ukf-example-k}\\
  y_t &= m_t+k_t+e_t,
    & e_t &\sim \calN(0,R). \label{eq:bf-ukf-example-y}
\end{align}
The state is $x_t=(m_t,k_t)$, but only the innovation
$\varepsilon_t$ is new stochastic input at time $t$.  Conditional on
$(m_{t-1},k_{t-1},\varepsilon_t)$, the coordinate $k_t$ is deterministic.

Suppose the previous filtering distribution is Gaussian:
\[
  x_{t-1}\mid y_{1:t-1}
  \approx \calN(\mu_{t-1},P_{t-1}),
  \qquad
  \mu_{t-1}=
  \begin{bmatrix}\bar m\\ \bar k\end{bmatrix}.
\]
The structural UKF should augment the previous state with the innovation.  This
step needs to be motivated carefully, because it is exactly where the structural
UKF departs from the naive full-state construction.  The point is not to pad the
state arbitrarily.  The point is to place sigma points on the variables that
actually generate the one-step uncertainty in the structural transition.  In
this example, the predictive law is generated by the joint uncertainty in the
lagged state $(m_{t-1},k_{t-1})$ and the current shock $\varepsilon_t$.  Once
those three quantities are fixed, both $m_t$ and $k_t$ are fixed by the
structural map.

That is why the augmented variable is
\begin{equation}
  a_t=(x_{t-1},\varepsilon_t),
  \label{eq:bf-ukf-augmented-variable}
\end{equation}
rather than an artificially padded post-transition full state.  Standard
additive-noise UKF presentations often augment with process noise for the same
reason: sigma points should live on the pre-transition uncertainty variables.
For a mixed structural model, however, that principle has a sharper consequence.
The shock $\varepsilon_t$ is a genuine new source of uncertainty, while $k_t$
is not.  If the augmentation omitted $\varepsilon_t$, the propagated sigma
points would miss current shock uncertainty.  If the augmentation instead added
a direct perturbation for $k_t$, the sigma points would describe a different
transition law from the one declared by the model.

So the structural UKF should augment the previous state with the innovation.
This places sigma points in the same pre-transition uncertainty variables that
define the predictive law, rather than in an artificially padded post-transition
full state:
\begin{equation}
  a_t =
  \begin{bmatrix}
    m_{t-1}\\ k_{t-1}\\ \varepsilon_t
  \end{bmatrix},
  \qquad
  a_t\mid y_{1:t-1}
  \approx
  \calN\left(
    \begin{bmatrix}\bar m\\ \bar k\\ 0\end{bmatrix},
    \begin{bmatrix}
      P_{t-1} & 0\\
      0 & 1
    \end{bmatrix}
  \right).
  \label{eq:bf-ukf-augmented-law}
\end{equation}
Let $\{a_t^{(j)},w_j^{(m)},w_j^{(c)}\}_{j=0}^{2n_a}$ be the usual unscented
points and weights for this $n_a=3$ dimensional augmented Gaussian
\citep{julier1997new}.  For every point
$a_t^{(j)}=(m_{t-1}^{(j)},k_{t-1}^{(j)},\varepsilon_t^{(j)})$, the structural
transition is evaluated as
\begin{align}
  m_t^{(j)}
    &= \rho m_{t-1}^{(j)}+\sigma\varepsilon_t^{(j)},\\
  k_t^{(j)}
    &= \phi k_{t-1}^{(j)}+\gamma\left(m_t^{(j)}\right)^2,\\
  x_t^{(j)}
    &= \begin{bmatrix}m_t^{(j)}\\ k_t^{(j)}\end{bmatrix}.
\end{align}
This is the crucial step: the filter never creates an independent sigma-point
perturbation for $k_t$.  It creates sigma points for the previous state and
the current innovation, then computes $k_t$ by the structural map
\eqref{eq:bf-ukf-example-k}.

The structural and naive full-state UKF constructions now differ at the level
that matters for the update.  The structural UKF forms sigma points for
$(x_{t-1},\varepsilon_t)$, propagates them through the structural transition,
and therefore computes moments of the target law declared by the model.  A
naive full-state UKF instead behaves as if there were an additive-noise model
of the form
\[
  x_t = \tilde f(x_{t-1}) + \tilde u_t,
  \qquad
  \tilde u_t \sim \calN(0,\tilde Q),
\]
with a full-state covariance $\tilde Q$ that assigns direct new variance to all
or most coordinates of $x_t$.  The three failure claims can now be derived in
the notation of this toy model.

\paragraph{Prediction-law failure.}
Under the structural model,
\[
  m_t = \rho m_{t-1} + \sigma \varepsilon_t,
  \qquad
  k_t = \phi k_{t-1} + \gamma m_t^2.
\]
So conditional on $(m_{t-1},k_{t-1})$, the one-step predictive law is induced by
a single scalar shock $\varepsilon_t$.  Eliminating the shock-driven update for
$m_t$ gives the structural identity
\begin{equation}
  k_t - \phi k_{t-1} - \gamma m_t^2 = 0.
  \label{eq:bf-ukf-structural-identity}
\end{equation}
Every propagated sigma point in the structural UKF satisfies this identity, so
its sigma-point cloud lies on the support of the structural predictive law.  A
naive full-state UKF instead propagates points
\[
  \tilde x_t^{(j)}=(\tilde m_t^{(j)},\tilde k_t^{(j)})
\]
from an additive full-state covariance model.  In general those points satisfy
\[
  \tilde k_t^{(j)} - \phi k_{t-1}^{(j)} - \gamma \bigl(\tilde m_t^{(j)}\bigr)^2
  \neq 0,
\]
so the naive sigma-point cloud approximates moments of a different one-step law.
That is the prediction-law failure.

\paragraph{Cross-covariance failure.}
For the observation equation $z_t=m_t+k_t$, the structural UKF computes
\[
  P_{xz,t}^{\mathrm{struct}}
  =
  \sum_j w_j^{(c)}
  \bigl(x_t^{(j)}-\hat x_{t|t-1}\bigr)
  \bigl(z_t^{(j)}-\hat z_t\bigr),
\]
with
\[
  z_t^{(j)} = m_t^{(j)} + k_t^{(j)},
  \qquad
  k_t^{(j)} = \phi k_{t-1}^{(j)} + \gamma \bigl(m_t^{(j)}\bigr)^2.
\]
Hence all variation in the second state coordinate enters through lagged-state
variation and the common shock channel already represented in $m_t^{(j)}$.
Under the naive full-state UKF,
\[
  P_{xz,t}^{\mathrm{naive}}
  =
  \sum_j w_j^{(c)}
  \bigl(\tilde x_t^{(j)}-\tilde x_{t|t-1}\bigr)
  \bigl(\tilde z_t^{(j)}-\tilde z_t\bigr),
\]
with
\[
  \tilde z_t^{(j)} = \tilde m_t^{(j)} + \tilde k_t^{(j)}.
\]
Because $\tilde k_t^{(j)}$ now contains a direct artificial perturbation channel,
the second component of $P_{xz,t}^{\mathrm{naive}}$ contains covariance terms
between the observation and that artificial variation.  Those terms do not exist
in the structural model, where $k_t$ is completed deterministically from the
same shock path that drives $m_t$.  That is the cross-covariance failure.

\paragraph{Update-law failure.}
The familiar UKF update equations,
\[
  K_t = P_{xz,t} S_t^{-1},
  \qquad
  \hat x_{t|t} = \hat x_{t|t-1} + K_t v_t,
  \qquad
  P_{t|t} = P_{t|t-1} - K_t S_t K_t^\top,
\]
do not fail algebraically.  They fail semantically when
$(\hat x_{t|t-1},S_t,P_{xz,t})$ are computed from the wrong predictive law.
Once the naive full-state approximation changes either $P_{xz,t}$ or $S_t$, it
changes the gain itself:
\[
  K_t^{\mathrm{naive}}
  =
  P_{xz,t}^{\mathrm{naive}} \left(S_t^{\mathrm{naive}}\right)^{-1}
  \neq
  P_{xz,t}^{\mathrm{struct}} \left(S_t^{\mathrm{struct}}\right)^{-1}
  = K_t^{\mathrm{struct}}
\]
in general.  The existing numerical comparison later in this example shows this
explicitly: adding artificial variance to $k_t$ changes the innovation variance
from $S_t=0.61217$ to $S_t=0.65217$, which then changes the approximate
log-likelihood contribution.  The same change propagates into the Kalman gain
and posterior update.  So the update law is not wrong as linear algebra; it is
correctly updating the wrong approximate model.

The predicted state moments are
\begin{align}
  \hat x_{t|t-1}
    &= \sum_j w_j^{(m)} x_t^{(j)},\\
  P_{xx,t}
    &= \sum_j w_j^{(c)}
       \left(x_t^{(j)}-\hat x_{t|t-1}\right)
       \left(x_t^{(j)}-\hat x_{t|t-1}\right)^\top .
\end{align}
For the observation equation \eqref{eq:bf-ukf-example-y}, define
\[
  z_t^{(j)}=m_t^{(j)}+k_t^{(j)}.
\]
Then the predicted observation, innovation variance, and cross covariance are
\begin{align}
  \hat z_t
    &= \sum_j w_j^{(m)} z_t^{(j)},\\
  S_t
    &= \sum_j w_j^{(c)}(z_t^{(j)}-\hat z_t)^2 + R,\\
  P_{xz,t}
    &= \sum_j w_j^{(c)}
       \left(x_t^{(j)}-\hat x_{t|t-1}\right)(z_t^{(j)}-\hat z_t).
\end{align}
The Gaussian measurement update is
\begin{align}
  v_t &= y_t-\hat z_t,\\
  K_t &= P_{xz,t}S_t^{-1},\\
  \hat x_{t|t} &= \hat x_{t|t-1}+K_tv_t,\\
  P_{t|t} &= P_{xx,t}-K_tS_tK_t^\top,
\end{align}
and the one-step approximate log likelihood contribution is
\[
  \ell_t
  =
  -\frac{1}{2}\left[
    \log(2\pi S_t)+\frac{v_t^2}{S_t}
  \right].
\]

\subsection{Reviewer Edge Case: Does \texorpdfstring{$\phi=0$}{phi = 0} Collapse the Example?}
\label{subsec:bf-phi-zero-collapse}

The model does not collapse just because $\phi=0$.  In that case
\[
  k_t=\gamma m_t^2,
  \qquad
  m_t=\rho m_{t-1}+\sigma\varepsilon_t,
\]
so $k_t$ is still random whenever $\gamma\neq0$ and the conditional law of
$m_t$ is nondegenerate.  What disappears is the inherited lagged-$k$ channel,
not the current-shock channel.  The support becomes the lower-dimensional
manifold
\[
  \{(m,k): k=\gamma m^2\},
\]
conditional on the parameters and the lagged state distribution.  It collapses
to a point only in degenerate cases, for example if $\gamma=0$ or if $m_t$ is
itself degenerate.  With the numerical calibration below but setting $\phi=0$,
the recomputed structural UKF moments are
\[
  \hat x_{t|t-1}
  =
  \begin{bmatrix}
    0\\0.11024
  \end{bmatrix},
  \qquad
  P_{xx,t}
  =
  \begin{bmatrix}
    0.27560 & 0\\
    0 & 0.04247
  \end{bmatrix},
\]
and the observation moments are
\[
  \hat z_t=0.11024,
  \qquad
  S_t=0.56807,
  \qquad
  P_{xz,t}=
  \begin{bmatrix}
    0.27560\\0.04247
  \end{bmatrix}.
\]
The $k$ variance is smaller because lagged $k_{t-1}$ no longer feeds into
$k_t$, but it is not zero: it is induced by the quadratic transformation of the
shock-driven $m_t$.

\subsection{Numerical Illustration}
\label{sec:bf-structural-ukf-numeric}

Use the parameters
\[
  \rho=0.8,\quad \sigma=0.5,\quad \phi=0.7,\quad
  \gamma=0.4,\quad R=0.25,
\]
and previous filtering moments
\[
  \mu_{t-1}=(0,0)^\top,
  \qquad
  P_{t-1}=\operatorname{diag}(0.04,0.09).
\]
Use the scaled unscented convention $\alpha=1$, $\beta=2$, and $\kappa=0$.
Since $n_a=3$, this gives $\lambda=\alpha^2(n_a+\kappa)-n_a=0$, central mean
weight $w_0^{(m)}=0$, central covariance weight $w_0^{(c)}=2$, and six
off-center mean/covariance weights equal to $1/6$.
The augmented covariance is
\[
  \operatorname{diag}(0.04,0.09,1),
\]
so $\sqrt{(n_a+\lambda)P_a}$ has diagonal entries
$0.34641$, $0.51962$, and $1.73205$.  The seven sigma points and their
deterministic structural completion are:
\[
\begin{array}{c|rrr|rrr}
  j
  & m_{t-1}^{(j)} & k_{t-1}^{(j)} & \varepsilon_t^{(j)}
  & m_t^{(j)} & k_t^{(j)} & z_t^{(j)}\\
  \hline
  0 & 0       & 0        & 0        & 0        & 0        & 0\\
  1 & 0.34641& 0        & 0        & 0.27713 & 0.03072 & 0.30785\\
  2 & 0      & 0.51962  & 0        & 0       & 0.36373 & 0.36373\\
  3 & 0      & 0        & 1.73205  & 0.86603 & 0.30000 & 1.16603\\
  4 & -0.34641& 0       & 0        & -0.27713& 0.03072 & -0.24641\\
  5 & 0      & -0.51962 & 0        & 0       & -0.36373& -0.36373\\
  6 & 0      & 0        & -1.73205 & -0.86603& 0.30000 & -0.56603
\end{array}
\]
For example, the third positive innovation point gives
$m_t=0.5(1.73205)=0.86603$ and
$k_t=0.4(0.86603)^2=0.30000$.  No row contains an independently sampled
``new'' shock for $k_t$.

Using the stated covariance weights, the structural
propagation gives
\[
  \hat x_{t|t-1}
  =
  \begin{bmatrix}
    0\\ 0.11024
  \end{bmatrix},
  \qquad
  P_{xx,t}
  =
  \begin{bmatrix}
    0.27560 & 0\\
    0 & 0.08657
  \end{bmatrix}.
\]
The predicted observation moments are
\[
  \hat z_t=0.11024,\qquad S_t=0.61217,\qquad
  P_{xz,t}=
  \begin{bmatrix}
    0.27560\\ 0.08657
  \end{bmatrix}.
\]
For an observation $y_t=0.30$,
\[
  v_t=0.18976,\qquad
  K_t=
  \begin{bmatrix}
    0.45020\\ 0.14141
  \end{bmatrix},
\]
and therefore
\[
  \hat x_{t|t}
  =
  \begin{bmatrix}
    0.08543\\ 0.13707
  \end{bmatrix},
  \qquad
  P_{t|t}
  =
  \begin{bmatrix}
    0.15152 & -0.03897\\
    -0.03897 & 0.07433
  \end{bmatrix}.
\]
The one-step log likelihood contribution is approximately
\[
  \ell_t=-0.70297.
\]
These numbers are illustrative calculations under the stated sigma-point
convention and should be read as a worked example rather than general
performance evidence.

\paragraph{Side-by-side prediction/update contrast.}
For this toy model, the structural UKF treats
\[
  a_t=(x_{t-1},\varepsilon_t)
\]
as the sigma-point object and computes propagated points by
\[
  m_t^{(j)} = \rho m_{t-1}^{(j)} + \sigma \varepsilon_t^{(j)},
  \qquad
  k_t^{(j)} = \phi k_{t-1}^{(j)} + \gamma \bigl(m_t^{(j)}\bigr)^2.
\]
It therefore forms
\[
  P_{xz,t}^{\mathrm{struct}}
  =
  \sum_j w_j^{(c)}
  \bigl(x_t^{(j)}-\hat x_{t|t-1}\bigr)
  \bigl(z_t^{(j)}-\hat z_t\bigr),
\]
where every $x_t^{(j)}$ lies on the support generated by the structural map.
A naive full-state UKF instead produces an artificial cloud
$\tilde x_t^{(j)}=(\tilde m_t^{(j)},\tilde k_t^{(j)})$ with
\[
  \tilde k_t^{(j)} \neq \phi k_{t-1}^{(j)} + \gamma \bigl(\tilde m_t^{(j)}\bigr)^2
  \text{ in general,}
\]
and therefore forms
\[
  P_{xz,t}^{\mathrm{naive}}
  =
  \sum_j w_j^{(c)}
  \bigl(\tilde x_t^{(j)}-\tilde x_{t|t-1}\bigr)
  \bigl(\tilde z_t^{(j)}-\tilde z_t\bigr),
\]
which mixes genuine shock-driven variation with artificial perturbations of the
deterministic-completion coordinate.

\paragraph{Structural identity test.}
A simple way to see the difference is to check whether propagated points satisfy
\[
  k_t - \phi k_{t-1} - \gamma m_t^2 = 0.
\]
The structural UKF satisfies this identity point by point for every propagated
sigma point.  A naive full-state UKF with direct perturbations to $k_t$ does
not.  That is the diagnostic place where the structural target and this naive
full-state approximation part company.

Now compare this with an off-manifold full-state approximation that adds an
independent artificial innovation variance $0.04$ directly to $k_t$ after the
structural map.  Nothing in
\eqref{eq:bf-ukf-example-m}--\eqref{eq:bf-ukf-example-y} justifies that extra
shock.  It changes the innovation variance from $S_t=0.61217$ to
$S_t=0.65217$ and changes the same observation's log likelihood contribution
from $-0.70297$ to $-0.73282$.  The numerical difference is small in this toy
calibration, but the conceptual difference is large: the latter filter is
approximating a different approximate state-space model relative to the
declared structural law.

To connect these numbers back to the three failures above, write the naive
full-state approximation as using
\[
  \tilde P_{xx,t}
  =
  \begin{bmatrix}
    0.27560 & 0\\
    0 & 0.12657
  \end{bmatrix},
\]
which differs from the structural covariance only by the artificial extra
variance in the $k_t$ coordinate.  Under the same observation equation
$z_t=m_t+k_t$, this gives
\[
  \tilde S_t = 0.65217,
  \qquad
  \tilde P_{xz,t} =
  \begin{bmatrix}
    0.27560\\ 0.12657
  \end{bmatrix},
\]
and therefore
\[
  \tilde K_t
  =
  \begin{bmatrix}
    0.42259\\ 0.19408
  \end{bmatrix},
  \qquad
  \tilde x_{t|t}
  =
  \begin{bmatrix}
    0.08019\\ 0.14707
  \end{bmatrix},
\]
for the same observation $y_t=0.30$.

These values make the three failures concrete:
\begin{enumerate}[label=(\roman*)]
  \item \textbf{Prediction-law failure:}
    the structural update uses
    $P_{xx,t}=\operatorname{diag}(0.27560,0.08657)$,
    while the naive full-state approximation uses
    $\tilde P_{xx,t}=\operatorname{diag}(0.27560,0.12657)$ because it has
    inserted artificial uncertainty directly into $k_t$.
  \item \textbf{Cross-covariance failure:}
    the structural update uses
    \[
      P_{xz,t}=
      \begin{bmatrix}
        0.27560\\ 0.08657
      \end{bmatrix},
    \]
    whereas the naive full-state approximation uses
    \[
      \tilde P_{xz,t}=
      \begin{bmatrix}
        0.27560\\ 0.12657
      \end{bmatrix}.
    \]
    The second component is larger only because the naive construction lets the
    observation load directly on artificial variation in $k_t$.
  \item \textbf{Update-law failure:}
    the structural gain and posterior mean are
    \[
      K_t=
      \begin{bmatrix}
        0.45020\\ 0.14141
      \end{bmatrix},
      \qquad
      \hat x_{t|t}=
      \begin{bmatrix}
        0.08543\\ 0.13707
      \end{bmatrix},
    \]
    whereas the naive full-state approximation gives
    \[
      \tilde K_t=
      \begin{bmatrix}
        0.42259\\ 0.19407
      \end{bmatrix},
      \qquad
      \tilde x_{t|t}=
      \begin{bmatrix}
        0.08019\\ 0.14707
      \end{bmatrix}.
    \]
    The update formula has not changed, but the gain and posterior correction do
    because they are now computed from the wrong predictive moments.
\end{enumerate}
So the numerical illustration does not merely show a slightly different
likelihood value.  It shows that once the naive approximation changes the
predictive covariance and cross covariance, the entire measurement update shifts
with it.

The lesson is not that UKF is wrong.  The lesson is that the UKF must be
applied to the correct structural transition.  Sigma points belong in the
declared stochastic integration space.  Endogenous deterministic coordinates
are completed by the model point by point.  In BayesFilter terms, this is a
design rule for mixed structural models and an approximation-label boundary for
any backend that instead enlarges the stochastic space.

\section[Degenerate Linear Transition Example]{Worked Example B: Degenerate Linear Transition with Nonlinear Measurement}
\label{sec:bf-structural-nonlinear-measurement}

The most practically relevant counterpart to the previous example is not a fully
stochastic linear-Gaussian state transition.  It is a degenerate linear
transition together with a nonlinear measurement equation.  That case is closer
to many DSGE and affine term-structure applications than a fully nonlinear
latent transition.  The key point is that the latent transition may still have a
structural split even when it is linear.  A nonlinear measurement equation does
not remove that split.

Consider the linear structural transition
\begin{align}
  m_t &= \rho m_{t-1} + \sigma\varepsilon_t,
    & \varepsilon_t &\sim \calN(0,1),
    \label{eq:bf-exp-affine-transition-m}\\
  k_t &= \phi k_{t-1} + \psi m_{t-1} + \chi m_t,
    \label{eq:bf-exp-affine-transition-k}\\
  x_t &= (m_t,k_t),
    \label{eq:bf-exp-affine-transition-state}
\end{align}
and the nonlinear exponential-affine measurement equation
\begin{align}
  y_t &= \exp(A^\top x_t + \mu) + e_t,
    & e_t &\sim \calN(0,R).
    \label{eq:bf-exp-affine-observation}
\end{align}
The latent transition is linear, but it is still structurally degenerate:
conditional on $(x_{t-1},\varepsilon_t)$, the endogenous coordinate $k_t$ is
completed deterministically from the same shock path that drives $m_t$.  So the
structural issue from the previous section does not disappear.  What changes is
where the nonlinearity appears.  The latent transition is linear, while the
measurement map is nonlinear.

This means there are now two layers to keep separate:
\begin{enumerate}[label=(\roman*)]
  \item the latent predictive law is still generated by the pre-transition
    uncertainty variables $(x_{t-1},\varepsilon_t)$, not by direct full-state
    perturbations of $x_t$;
  \item even after the latent cloud is propagated correctly, the UKF must still
    approximate nonlinear observation moments and cross covariances for the map
    $x_t \mapsto \exp(A^\top x_t + \mu)$.
\end{enumerate}
The first issue is structural-modeling correctness.  The second is observation-
side quadrature quality.

\subsection{Worked structural derivation for the degenerate linear-transition case}
\label{sec:bf-exp-affine-formal}

\paragraph{Accepted assumptions.}
For this section, treat the following as explicit modeling assumptions:
\begin{enumerate}[label=(A\arabic*)]
  \item the innovation is conditionally independent of the lagged filtering law,
    so that
    \begin{equation}
      p(x_{t-1},\varepsilon_t\mid y_{1:t-1})
      =
      p(x_{t-1}\mid y_{1:t-1})p(\varepsilon_t);
      \label{ass:bf-exp-affine-factorization}
    \end{equation}
  \item the latent map, observation map, and composed map defined below are
    measurable and have the finite moments needed for the stated expectations;
  \item whenever we refer to the observation-side-only error regime, we assume
    that the latent predictive law from
    \eqref{eq:bf-exp-affine-pushforward} is preserved exactly;
  \item no invertibility of $T_k$ is required: it is enough that $k_t$ is a
    well-defined measurable deterministic completion of
    $(k_{t-1},m_{t-1},m_t)$;
  \item if one wants to replace the innovation variable $\varepsilon_t$ by the
    current exogenous state $m_t$ as the sigma-point variable, then for fixed
    $m_{t-1}$ the map
    \begin{equation}
      \varepsilon_t \mapsto T_m(m_{t-1},\varepsilon_t)
      \label{ass:bf-exp-affine-Tm-reparameterization}
    \end{equation}
    should be injective, or invertible onto its image with measurable inverse.
    This extra assumption is only needed for that reparameterization step, not
    for the pushforward identities proved below.
\end{enumerate}

\paragraph{Latent predictive pushforward.}
\label{lem:bf-exp-affine-latent-pushforward}
Define the degenerate linear structural map
\begin{equation}
  F_{\mathrm{lin}}(x_{t-1},\varepsilon_t)
  =
  \Bigl(
    \rho m_{t-1}+\sigma\varepsilon_t,
    \phi k_{t-1}+\psi m_{t-1}+\chi(\rho m_{t-1}+\sigma\varepsilon_t)
  \Bigr).
  \label{eq:bf-exp-affine-latent-map}
\end{equation}
Then for any bounded test function $\varphi$ on the latent-state space,
\begin{equation}
  \E\!\left[\varphi(x_t)\mid y_{1:t-1}\right]
  =
  \iint
  \varphi\!\left(F_{\mathrm{lin}}(x_{t-1},\varepsilon_t)\right)
  p(x_{t-1}\mid y_{1:t-1})p(\varepsilon_t)
  \, d\varepsilon_t \, dx_{t-1}.
  \label{eq:bf-exp-affine-pushforward}
\end{equation}
Hence the latent predictive law is the pushforward of the joint law of
$(x_{t-1},\varepsilon_t)$ through $F_{\mathrm{lin}}$.

\emph{Derivation.}
Under the state transition
\eqref{eq:bf-exp-affine-transition-m}--\eqref{eq:bf-exp-affine-transition-state}
and the factorization assumption
\eqref{ass:bf-exp-affine-factorization}, the conditional expectation of any
bounded test function $\varphi(x_t)$ is obtained by substituting
$x_t=F_{\mathrm{lin}}(x_{t-1},\varepsilon_t)$ and integrating over the joint law
of $(x_{t-1},\varepsilon_t)$.  This gives
\eqref{eq:bf-exp-affine-pushforward}, which is exactly the pushforward identity.

\paragraph{Deterministic-completion identity.}
\label{lem:bf-exp-affine-deterministic-completion}
Under \eqref{eq:bf-exp-affine-transition-m}--\eqref{eq:bf-exp-affine-transition-k},
all admissible latent states satisfy
\begin{equation}
  k_t - \phi k_{t-1} - \psi m_{t-1} - \chi m_t = 0.
  \label{eq:bf-exp-affine-structural-identity}
\end{equation}
Therefore, once $(x_{t-1},\varepsilon_t)$ is fixed, both $m_t$ and $k_t$ are
fixed; there is no independent new perturbation attached to $k_t$ in the latent
transition.

\emph{Derivation.}
Substitute \eqref{eq:bf-exp-affine-transition-m} into
\eqref{eq:bf-exp-affine-transition-k} and rearrange terms.  The resulting
identity is exactly \eqref{eq:bf-exp-affine-structural-identity}.

\paragraph{Observation predictive pushforward.}
\label{lem:bf-exp-affine-observation-pushforward}
Define the nonlinear noise-free observation map
\begin{equation}
  h(x_t)=\exp(A^\top x_t+\mu).
  \label{eq:bf-exp-affine-observation-map}
\end{equation}
and the composed map
\begin{equation}
  G(x_{t-1},\varepsilon_t)
  =
  h\!\left(F_{\mathrm{lin}}(x_{t-1},\varepsilon_t)\right).
  \label{eq:bf-exp-affine-composed-map}
\end{equation}
Then for any bounded test function $\psi_0$ on the noise-free observation
space,
\begin{equation}
  \E\!\left[\psi_0\!\bigl(h(x_t)\bigr)\mid y_{1:t-1}\right]
  =
  \iint
  \psi_0\!\left(G(x_{t-1},\varepsilon_t)\right)
  p(x_{t-1}\mid y_{1:t-1})p(\varepsilon_t)
  \, d\varepsilon_t \, dx_{t-1}.
  \label{eq:bf-exp-affine-observation-pushforward}
\end{equation}
Hence the noise-free observation predictive law is the pushforward of the same
pre-transition uncertainty through the composed map.

\emph{Derivation.}
Compose the latent pushforward identity
\eqref{eq:bf-exp-affine-pushforward} with the measurable observation map
\eqref{eq:bf-exp-affine-observation-map}.  This yields
\eqref{eq:bf-exp-affine-observation-pushforward}.

\paragraph{Naive full-state perturbation changes the latent law.}
\label{lem:bf-exp-affine-naive-full-state-support-violation}
If a naive full-state UKF introduces a direct perturbation to the endogenous
coordinate $k_t$, then its propagated latent sigma-point cloud does not, in
general, satisfy \eqref{eq:bf-exp-affine-structural-identity} pointwise.
Therefore it does not target the latent predictive law characterized by
\eqref{eq:bf-exp-affine-pushforward}.

\emph{Diagnostic check.}
By the deterministic-completion identity above, every admissible
latent state under the structural model satisfies
\eqref{eq:bf-exp-affine-structural-identity}.  A naive full-state UKF that adds
an independent perturbation to $k_t$ creates propagated points whose $k_t$
values are not determined solely by $(x_{t-1},\varepsilon_t)$, so the identity
fails in general.  Hence the resulting sigma-point cloud cannot represent the
same latent predictive law.

\paragraph{Failure propagation and observation-side-only error.}
\label{prop:bf-exp-affine-law-mismatch-propagation}
Under the preceding worked identities:
\begin{enumerate}[label=(\roman*)]
  \item if a naive full-state UKF directly perturbs $k_t$, then
    prediction-law failure, cross-covariance failure, and update-law failure all
    still apply;
  \item if the latent predictive law from
    \eqref{eq:bf-exp-affine-pushforward} is preserved exactly and only the
    transformed observation moments from
    \eqref{eq:bf-exp-affine-observation-pushforward} are approximated by
    quadrature, then the error is not structural-model failure; it is only
    observation-side quadrature error.
\end{enumerate}

\emph{Reason.}
Part~(i) follows because the naive construction changes the latent predictive
law, and the nonlinear observation law is the pushforward of that same latent
law.  Once the latent law changes, the observation moments, state-observation
cross covariance, and therefore the UKF update objects also change.  Part~(ii)
follows from Assumption (A3): if the latent law is preserved and only the
nonlinear observation moments are approximated, then the remaining error arises
only in the observation-side quadrature.

The formal statements above separate three layers that are easy to blur in
practice: accepted modeling assumptions, exact pushforward identities, and the
interpretation of failure modes.  The filtered state is still $x_t$, but the
latent predictive law is generated by the pre-transition uncertainty variables
$(x_{t-1},\varepsilon_t)$.  The nonlinear measurement equation adds another
approximation layer, but it does not remove the need to respect the degenerate
latent transition.  Assumption~(A5) is stronger than the pushforward arguments
need: it is included only to justify replacing $\varepsilon_t$ by $m_t$ as an
equivalent sigma-point variable when that reparameterization is desired.

\subsection{Worked Degenerate Linear-Transition Example}
\label{sec:bf-exp-affine-example}

Use
\[
  \rho = 0.8,
  \qquad
  \sigma = 0.5,
  \qquad
  \phi = 0.7,
  \qquad
  \psi = 0.3,
  \qquad
  \chi = 0.4,
\]
with previous filtering moments
\[
  \mu_{t-1}=(0,0)^\top,
  \qquad
  P_{t-1}=\operatorname{diag}(0.04,0.09),
\]
and measurement parameters
\[
  A=
  \begin{bmatrix}
    1.1\\ 0.6
  \end{bmatrix},
  \qquad
  \mu = 0.10,
  \qquad
  R = 0.04.
\]
Because the latent transition remains structurally degenerate, the correct
sigma-point object is still
\[
  a_t=(x_{t-1},\varepsilon_t),
\]
not a naively perturbed full post-transition state.  Using the same scaled
unscented construction as in the previous example, the propagated latent points
are obtained from
\begin{align}
  m_t^{(j)} &= \rho m_{t-1}^{(j)} + \sigma\varepsilon_t^{(j)},\\
  k_t^{(j)} &= \phi k_{t-1}^{(j)} + \psi m_{t-1}^{(j)} + \chi m_t^{(j)}.
\end{align}
For each propagated point, define the nonlinear noise-free observation
\[
  z_t^{(j)} = \exp\!\left(1.1 m_t^{(j)} + 0.6 k_t^{(j)} + 0.10\right).
\]
The same seven sigma points used earlier now produce the following latent and
measurement values:
\[
\begin{array}{c|rrr|rrr}
  j
  & m_{t-1}^{(j)} & k_{t-1}^{(j)} & \varepsilon_t^{(j)}
  & m_t^{(j)} & k_t^{(j)} & z_t^{(j)}\\
  \hline
  0 & 0 & 0 & 0 & 0 & 0 & 1.10517\\
  1 & 0.34641 & 0 & 0 & 0.27713 & 0.21477 & 1.70524\\
  2 & 0 & 0.51962 & 0 & 0 & 0.36373 & 1.37470\\
  3 & 0 & 0 & 1.73205 & 0.86603 & 0.34641 & 3.52709\\
  4 & -0.34641 & 0 & 0 & -0.27713 & -0.21477 & 0.71626\\
  5 & 0 & -0.51962 & 0 & 0 & -0.36373 & 0.88848\\
  6 & 0 & 0 & -1.73205 & -0.86603 & -0.34641 & 0.34629
\end{array}
\]
The predicted latent moments are
\[
  \hat x_{t|t-1}
  =
  \begin{bmatrix}
    0\\ 0
  \end{bmatrix},
  \qquad
  P_{xx,t}
  =
  \begin{bmatrix}
    0.27560 & 0.11984\\
    0.11984 & 0.09948
  \end{bmatrix}.
\]
The nonlinear predicted observation moments are
\[
  \hat z_t = 1.42635,
  \qquad
  S_t = 1.32191,
  \qquad
  P_{xz,t}
  =
  \begin{bmatrix}
    0.50479\\ 0.24852
  \end{bmatrix}.
\]
So the structural UKF gain is
\[
  K_t
  =
  \begin{bmatrix}
    0.38186\\ 0.18800
  \end{bmatrix}.
\]
If the realized observation is $y_t=1.50$, then
\[
  v_t = 0.07365,
  \qquad
  \hat x_{t|t}
  =
  \begin{bmatrix}
    0.02813\\ 0.01385
  \end{bmatrix}.
\]
As in the previous example, these numbers are illustrative calculations under
the stated sigma-point convention.  They demonstrate the law and moment objects
being compared; they are not standalone performance evidence for a production
filter.

Now compare this with a naive full-state approximation that adds an artificial
latent perturbation variance $0.04$ directly to the deterministic-completion
coordinate $k_t$.  To keep the comparison reproducible, take the naive rule to
form a four-dimensional sigma-point variable
\[
  \tilde a_t=(x_{t-1},\varepsilon_t,\delta_t),
  \qquad
  \delta_t\sim\calN(0,0.04),
\]
and to propagate
\[
  \tilde k_t
  =
  \phi k_{t-1}+\psi m_{t-1}+\chi m_t+\delta_t.
\]
That artificial perturbation changes the latent predictive covariance to
\[
  \tilde P_{xx,t}
  =
  \begin{bmatrix}
    0.27560 & 0.11984\\
    0.11984 & 0.13948
  \end{bmatrix},
\]
while leaving the mean unchanged.  Under the same nonlinear
measurement map, the corresponding naive moment approximation is
\[
  \tilde z_t = 1.44483,
  \qquad
  \tilde S_t = 1.57838,
  \qquad
  \tilde P_{xz,t}
  =
  \begin{bmatrix}
    0.53757\\ 0.28867
  \end{bmatrix}.
\]
Hence the naive full-state gain becomes
\[
  \tilde K_t
  =
  \begin{bmatrix}
    0.34058\\ 0.18289
  \end{bmatrix},
\]
and for the same observation $y_t=1.50$ the naive posterior mean update is
\[
  \tilde x_{t|t}
  =
  \begin{bmatrix}
    0.01879\\ 0.01009
  \end{bmatrix}.
\]
These numbers make the three failures concrete in this linear-transition /
nonlinear-measurement setting:
\begin{enumerate}[label=(\roman*)]
  \item \textbf{Prediction-law failure:}
    the structural latent law uses
    \[
      P_{xx,t}
      =
      \begin{bmatrix}
        0.27560 & 0.11984\\
        0.11984 & 0.09948
      \end{bmatrix},
    \]
    whereas the naive full-state approximation uses
    \[
      \tilde P_{xx,t}
      =
      \begin{bmatrix}
        0.27560 & 0.11984\\
        0.11984 & 0.13948
      \end{bmatrix},
    \]
    because it has inserted artificial uncertainty directly into $k_t$.
  \item \textbf{Cross-covariance failure:}
    the structural update uses
    \[
      P_{xz,t}
      =
      \begin{bmatrix}
        0.50479\\ 0.24852
      \end{bmatrix},
    \]
    whereas the naive approximation uses
    \[
      \tilde P_{xz,t}
      =
      \begin{bmatrix}
        0.53757\\ 0.28867
      \end{bmatrix}.
    \]
    The extra direct perturbation of $k_t$ changes how the nonlinear
    observation loads on the second state coordinate.
  \item \textbf{Update-law failure:}
    the structural gain and posterior mean are
    \[
      K_t
      =
      \begin{bmatrix}
        0.38186\\ 0.18800
      \end{bmatrix},
      \qquad
      \hat x_{t|t}
      =
      \begin{bmatrix}
        0.02813\\ 0.01385
      \end{bmatrix},
    \]
    whereas the naive full-state approximation gives
    \[
      \tilde K_t
      =
      \begin{bmatrix}
        0.34058\\ 0.18289
      \end{bmatrix},
      \qquad
      \tilde x_{t|t}
      =
      \begin{bmatrix}
        0.01879\\ 0.01009
      \end{bmatrix}.
    \]
    The update formula itself is unchanged, but the predicted moments and cross
    covariance are not, so the gain and posterior correction are computed for a
    different approximate model.
\end{enumerate}

\subsection{What Is Similar to the Structural-Deterministic Case?}
\label{sec:bf-exp-affine-similarities}

The similarity is that the structural latent transition is still generated by
$(x_{t-1},\varepsilon_t)$, not by independent perturbations of every coordinate
of $x_t$.  If a naive full-state UKF directly perturbs $k_t$, it again changes
the latent predictive law before the observation map is even evaluated.  So the
structural warning remains active even though the latent transition is linear.

The nonlinear observation adds a second approximation layer.  Once the latent
cloud is propagated correctly, the UKF still has to approximate
\[
  \hat y_t,
  \qquad
  S_t,
  \qquad
  P_{xz,t}
\]
for the map $x_t \mapsto \exp(A^\top x_t + \mu)$.  Poor quadrature here still
changes the gain and posterior update, just as in other nonlinear Gaussian
filters.

\subsection{What Is Different from the Earlier Nonlinear-State Example?}
\label{sec:bf-exp-affine-differences}

In the earlier toy example, the latent state equation itself was nonlinear, so
the UKF had to approximate both the latent pushforward through the nonlinear
state map and the resulting observation moments.  Here the latent transition is
linear, but structurally degenerate.  That removes one source of approximation
error but not the structural state-partition issue.  The UKF still needs the
correct pre-transition uncertainty variable because $k_t$ remains a
deterministic-completion coordinate.

So this example isolates a different practical regime:
\begin{enumerate}[label=(\roman*)]
  \item the latent transition is linear but not full-state stochastic;
  \item the nonlinear burden sits in the measurement equation;
  \item the UKF can still fail in two distinct ways:
    \begin{itemize}
      \item structurally, if it perturbs $k_t$ directly and therefore changes the
        latent predictive law;
      \item numerically, if it uses the correct latent construction but poorly
        approximates the nonlinear observation moments.
    \end{itemize}
\end{enumerate}
This is exactly why the structural deterministic-dynamics chapter should not be
read as applying only to nonlinear state transitions.  The same latent
state-partition discipline remains necessary when the state equation is linear
but degenerate and the observation map is nonlinear.

\section[Structural Correctness Boundary]{What Structural Correctness Does and Does Not Guarantee}
\label{sec:bf-structural-correctness-boundary}

Structural correctness guarantees that the filter is approximating the intended latent predictive law rather than an altered full-state approximation.  It does not, by itself, guarantee that the resulting Gaussian closure is exact, that nonlinear observation moments are approximated with negligible error, or that the derivative path is stable enough for HMC.  In BayesFilter terms, structural correctness is a law-specification requirement; it is not a blanket certificate of moment accuracy, derivative accuracy, or production readiness.

In particular, structural correctness does not solve four other problems that must be evaluated separately:
\begin{enumerate}[label=(\roman*)]
  \item it does not make the Gaussian sigma-point closure exact for a genuinely nonlinear transformation;
  \item it does not remove observation-side quadrature error in models such as
    \eqref{eq:bf-exp-affine-observation};
  \item it does not resolve numerical covariance-factor issues when square roots,
    SVD factors, or PSD projections are used;
  \item it does not guarantee a stable value-gradient path for HMC, especially
    when spectral decompositions or fallback branches are differentiated.
\end{enumerate}
Thus the right mental model is layered: first make sure the latent law is correct, then measure the quality of the nonlinear approximation, then audit the numerical and derivative path.

\section[Degeneracy and SVD]{Structural Degeneracy, Numerical Degeneracy, and SVD}
\label{sec:bf-structural-versus-numerical-degeneracy}

The literature and the source-project audit motivate a distinction that should
be made explicit:
\begin{center}
\begin{tabular}{p{0.22\linewidth}p{0.30\linewidth}p{0.34\linewidth}}
\toprule
\textbf{Issue} & \textbf{What it means} & \textbf{What fixes it}\\
\midrule
Structural degeneracy & Some latent coordinates are deterministic-completion variables under the model law & Respect the declared stochastic integration space and complete the deterministic block pointwise\\
Numerical degeneracy & A covariance or innovation matrix is singular or ill-conditioned in computation & Use a stable solve, square-root, or spectral factorization appropriate to the declared law\\
SVD/square-root representation & A numerical backend chooses a robust factorization of a covariance object & This may improve value-side robustness, but it does not by itself correct a structurally wrong latent law\\
Derivative/spectral risk & Raw autodiff through spectral or fallback branches may be unstable & Requires derivative audits, stress tests, and possibly custom gradients\\
\bottomrule
\end{tabular}
\end{center}
This is why an SVD sigma-point filter does not automatically solve the problem
discussed in this chapter.  It may solve a factorization problem while leaving
the latent-law specification problem untouched.  Conversely, a structural SVD
sigma-point filter is a sensible backend when it factors the covariance of the
declared pre-transition variable, propagates those points through
$F_\theta$, and records the usual value and derivative diagnostics from
Chapters~\ref{ch:bf-square-root-sigma-point} and
\ref{ch:bf-svd-sigma-point}.  The question is therefore not
``SVD or structural UKF?''  The correct question is:
\[
  \text{Which law is being factored and propagated?}
\]

\begin{proposition}[Degenerate covariance and formal sigma-point accuracy]
\label{prop:bf-degenerate-covariance-formal-ut-accuracy}
Let $A_\delta=\mu+\delta Z$ be a possibly degenerate Gaussian random
variable in $\R^L$, with $\E[Z]=0$ and
$\Cov(Z)=\Sigma\succeq0$.  Thus
$\Cov(A_\delta)=P_\delta=\delta^2\Sigma$ may be singular.  Let
$\{a_\delta^{(j)},w_j^{(m)},w_j^{(c)}\}$ be any sigma-point rule whose
deviations $\xi_j=a_\delta^{(j)}-\mu=\delta\zeta_j$ have fixed standardized
points $\zeta_j$ and satisfy
\[
  \sum_j w_j^{(m)}\xi_j=0,\qquad
  \sum_j w_j^{(m)}\xi_j\xi_j^\top=P_\delta,\qquad
  \sum_j w_j^{(c)}\xi_j\xi_j^\top=P_\delta .
\]
If $G:\R^L\to\R^d$ is three-times continuously differentiable and its
third-order Taylor remainders around $\mu$ have expected size $O(\delta^3)$
under the Gaussian law and weighted size $O(\delta^3)$ under the sigma-point
cloud, then the sigma-point transformed mean and covariance have the same local
order as in the nonsingular case:
\begin{align}
  \E[G_r(A_\delta)]
    &=
    G_r(\mu)+\frac{1}{2}\tr(H_rP_\delta)+O(\delta^3),\\
  \hat g_r^{\mathrm{SP}}
    &=
    G_r(\mu)+\frac{1}{2}\tr(H_rP_\delta)+O(\delta^3),\\
  \Cov(G(A_\delta))
    &=
    JP_\delta J^\top+O(\delta^3),\\
  P_G^{\mathrm{SP}}
    &=
    JP_\delta J^\top+O(\delta^3),
\end{align}
where $J$ is the Jacobian of $G$ at $\mu$ and $H_r$ is the Hessian of
the $r$th component.  Positive definiteness of $P_\delta$ is not used.
Consequently, abstracting from finite arithmetic and factor-construction
errors, rank deficiency of the state-noise covariance does not by itself lower
the formal local sigma-point moment accuracy for the law being propagated.
\end{proposition}

\begin{proof}
The usual Taylor proof only uses the first two moments of the input variable
and the matching first two weighted moments of the sigma-point cloud.  For a
component $r$,
\[
  G_r(\mu+\Delta)
  =
  G_r(\mu)+J_r\Delta+\frac{1}{2}\Delta^\top H_r\Delta
  +O(\|\Delta\|^3).
\]
Taking expectations with $\Delta=A_\delta-\mu$ gives the first expansion
because $\E[\Delta]=0$ and $\E[\Delta\Delta^\top]=P_\delta$.  Taking the
weighted sum with $\Delta=\xi_j$ gives the sigma-point mean expansion by the
same identities.  For covariance, the leading centered term is
$J\Delta$, whose second moment is $JP_\delta J^\top$ both for the true
Gaussian law and for the sigma-point cloud.  The quadratic remainder terms
enter at order $O(\delta^3)$ or smaller.  No inverse, Cholesky pivot, or
full-rank support argument appears in this proof, so the conclusion is
unchanged when $P_\delta$ is only positive semidefinite.
\end{proof}

Proposition~\ref{prop:bf-degenerate-covariance-formal-ut-accuracy} verifies
the reviewer's claim in its precise value-side form.  If a model is written as
a full-state transition with a degenerate Gaussian innovation, and if that
transition induces exactly the same conditional law as the structural
pushforward in Proposition~\ref{prop:bf-structural-ukf-pushforward}, then an
exact SVD sigma-point rule for that degenerate law has the same formal local
moment order as the corresponding nondegenerate rule.  Degeneracy is not, by
itself, a mathematical objection to sigma-point filtering.
This is a formal value-side statement abstracting from finite arithmetic,
factor-construction choices, and derivative-branch effects.  Those
implementation issues still require the diagnostics below.

That statement is an equivalence result, not a recommendation to hide the
structural split inside a large singular covariance matrix.  To make the
comparison precise, let the structural route use
\[
  A_t=(x_{t-1},\varepsilon_t),\qquad
  x_t=F_\theta(A_t),\qquad
  \pi_t^x=(F_\theta)_\#\pi_t^a ,
\]
where $\pi_t^a=\pi_{t-1|t-1}^x\otimes\nu_\theta^\varepsilon$.  Let a collapsed
full-state SVD route instead use an augmented variable
\[
  \widetilde A_t=(x_{t-1},\widetilde u_t),\qquad
  \widetilde u_t\sim\calN(0,\widetilde Q_t),\qquad
  \widetilde x_t=\widetilde F_\theta(\widetilde A_t),
\]
with $\widetilde Q_t=\widetilde C_t\widetilde C_t^\top$ obtained by an SVD or
spectral rule.  The collapsed route is acceptable only when the following
mathematical diagnostics are satisfied.

\begin{enumerate}[label=(\roman*)]
  \item \textbf{Law declaration is pushforward equality.}  The two
    declarations are equivalent only if
    \[
      (\widetilde F_\theta)_\#
      \left(\pi_{t-1|t-1}^x\otimes\calN(0,\widetilde Q_t)\right)
      =
      (F_\theta)_\#
      \left(\pi_{t-1|t-1}^x\otimes\nu_\theta^\varepsilon\right),
      \label{eq:bf-collapsed-law-equivalence}
    \]
    or, equivalently, if
    $\E[\varphi(\widetilde x_t)\mid y_{1:t-1}]
    =\E[\varphi(x_t)\mid y_{1:t-1}]$ for every bounded test function
    $\varphi$.  Matching a covariance matrix is only a necessary moment check,
    not this equality of laws.

    \emph{Example.}  If $m_t=\varepsilon_t$ and $k_t=\varepsilon_t^2$, then the
    structural conditional law is supported on the parabola
    $\{(m,k):k=m^2\}$.  A collapsed additive Gaussian transition
    $\widetilde x_t=\mu+\widetilde u_t$ with any covariance
    $\widetilde Q_t$ has affine Gaussian support.  It cannot equal the
    parabolic pushforward unless $\varepsilon_t$ is degenerate.  A singular
    covariance matrix by itself therefore does not encode the nonlinear timing
    contract.

  \item \textbf{Sigma-point support is a constraint residual.}  Let
    $c_\theta(a,x)=0$ collect the deterministic identities implied by the
    structural transition, for example
    $c_\theta(A_t,x_t)=x_t-F_\theta(A_t)$.  A structural sigma-point cloud
    satisfies
    \[
      \Delta_{\mathrm{id}}^{\mathrm{str}}
      =
      \max_j \|c_\theta(a_t^{(j)},F_\theta(a_t^{(j)}))\|
      =0.
    \]
    The collapsed cloud must satisfy the corresponding check
    \[
      \Delta_{\mathrm{id}}^{\mathrm{col}}
      =
      \max_j
      \|c_\theta(\widetilde a_t^{(j)},
        \widetilde F_\theta(\widetilde a_t^{(j)}))\|
      =0,
      \label{eq:bf-collapsed-identity-residual}
    \]
    after translating its variables into the structural coordinates.  If this
    residual is not defined because the collapsed adapter no longer exposes the
    relevant shock or timing variables, the adapter cannot demonstrate
    structural equivalence.

    \emph{Example.}  For the linear deterministic identity
    $x_t=(m_t,k_t)=(\varepsilon_t,\varepsilon_t)$, the support condition is
    $c(m,k)=k-m=0$.  The exact rank-one covariance
    $\widetilde Q=\begin{psmallmatrix}1&1\\1&1\end{psmallmatrix}$ has an SVD
    factor whose points lie on $k=m$.  After a nugget,
    $\widetilde Q_\eta=\widetilde Q+\eta^2I$, sigma points have normal
    coordinate displacement $\sqrt{L+\lambda}\,\eta$ in the direction
    $(1,-1)/\sqrt2$, so $\Delta_{\mathrm{id}}^{\mathrm{col}}>0$.

  \item \textbf{Dimension and cost are point-count identities.}  For a
    symmetric $2L+1$ rule, the number of model evaluations is determined by the
    dimension $L$ of the factored variable.  If the previous state has dimension
    $n$ and the declared shock has dimension $q$, the structural augmented rule
    uses
    \[
      N_{\mathrm{str}}=2(n+q)+1.
    \]
    A full-state additive-noise route that augments by an $n$-dimensional
    degenerate noise vector uses
    \[
      N_{\mathrm{col}}=2(n+n)+1,
      \qquad
      N_{\mathrm{col}}-N_{\mathrm{str}}=2(n-q).
      \label{eq:bf-collapsed-point-count}
    \]
    This overhead is pure bookkeeping if the extra $n-q$ directions are
    zero-noise directions.

    \emph{Example.}  With $n=100$ state coordinates and $q=10$ genuine shocks,
    the structural augmented rule uses $221$ points, while the full-state
    degenerate-noise augmentation uses $401$ points.  In the two-coordinate toy
    covariance $\diag(1,0)$ treated as a two-dimensional factor, the zero
    column produces two sigma points equal to the mean; they add evaluations
    and alter the weight bookkeeping without adding stochastic information.

  \item \textbf{Finite arithmetic is normal-support leakage.}  Let $M_t$ be a
    matrix whose columns span the exact linear support of a degenerate
    Gaussian factor, and let $N_t$ span the normal space
    $(\operatorname{col}M_t)^\perp$.  Exact support preservation requires
    \[
      N_t^\top \widetilde C_t=0
      \quad\Longleftrightarrow\quad
      N_t^\top\widetilde Q_tN_t=0 .
      \label{eq:bf-collapsed-normal-support}
    \]
    A nugget, PSD projection, rank threshold, or asymmetric reconstruction
    should therefore be reported through
    \[
      \ell_{\perp,t}
      =
      \|N_t^\top\widetilde Q_tN_t\|_{\mathrm{op}}
      \quad\text{or}\quad
      \max_j\|N_t^\top(\widetilde x_t^{(j)}-\bar x_t)\|.
      \label{eq:bf-collapsed-support-leakage}
    \]

    \emph{Example.}  In the identity example $k=m$, the normal vector is
    $N=(1,-1)^\top/\sqrt2$.  For
    $\widetilde Q_\eta=
    \begin{psmallmatrix}1&1\\1&1\end{psmallmatrix}+\eta^2I$,
    \[
      N^\top\widetilde Q_\eta N=\eta^2.
    \]
    Hence the nominally deterministic relation has positive artificial
    variance $\eta^2$ in finite arithmetic.  This is not a harmless
    representation of the same law; it is a regularized law unless explicitly
    removed before propagation.

  \item \textbf{Gradient risk is controlled by spectral gaps.}  If
    $\widetilde C_t(\theta)$ is produced from an eigendecomposition or SVD,
    derivatives of the spectral vectors contain gap denominators.  For a
    symmetric eigendecomposition, schematically,
    \[
      d u_i
      =
      \sum_{j\ne i}
      u_j
      \frac{u_j^\top(d\widetilde Q_t)u_i}{\lambda_i-\lambda_j}
      +\text{terms along }u_i .
      \label{eq:bf-collapsed-spectral-gap-gradient}
    \]
    Therefore the gradient audit must track
    \[
      g_t^{\min}
      =
      \min_{i\ne j}|\lambda_i(\widetilde Q_t)-\lambda_j(\widetilde Q_t)|,
      \qquad
      \|\partial_\theta\widetilde C_t\|
      \lesssim
      \frac{\|\partial_\theta\widetilde Q_t\|}{g_t^{\min}}
    \]
    away from rank and ordering branches.  Structural completion does not
    remove factor-gradient issues, but it avoids adding spectral choices for
    artificial zero-noise coordinates.

    \emph{Example.}  Let
    $\widetilde Q(\theta)=
    \begin{psmallmatrix}1&\theta\\ \theta&1\end{psmallmatrix}$.
    The eigenvalues are $1+\theta$ and $1-\theta$, with gap $2|\theta|$.
    At $\theta=0$ the covariance is perfectly smooth but the ordered
    eigenspace is not unique; raw autodiff through the spectral factor can be
    discontinuous or unstable even though the value-side covariance is benign.

  \item \textbf{Diagnostics must separate model law from regularization.}  A
    collapsed route should report both the identity residual
    \eqref{eq:bf-collapsed-identity-residual} and the normal leakage
    \eqref{eq:bf-collapsed-support-leakage}.  Otherwise a finite likelihood can
    be produced by a different state-space model.  For a scalar observation
    that loads on a deterministic coordinate, this difference is explicit:
    if the structural law is
    \[
      k_t=0,\qquad y_t=k_t+e_t,\qquad e_t\sim\calN(0,r),
    \]
    then $y_t\sim\calN(0,r)$.  If a collapsed covariance injects artificial
    variance $\eta^2$ into $k_t$, then the implemented likelihood uses
    $y_t\sim\calN(0,r+\eta^2)$, and the one-period log-likelihood change is
    \[
      \Delta\ell_t
      =
      -\frac12
      \left[
        \log\frac{r+\eta^2}{r}
        +
        y_t^2
        \left(\frac{1}{r+\eta^2}-\frac{1}{r}\right)
      \right].
      \label{eq:bf-collapsed-nugget-likelihood-change}
    \]
    A finite value of this likelihood is therefore not evidence that the
    structural transition was respected; it may only show that the nugget made
    the altered model numerically evaluable.
\end{enumerate}

These diagnostics matter for the UKF because the UKF update is only as correct
as the sigma-point cloud used to build its predictive moments.  For a
noise-free observation map $h_\theta$ and observation noise covariance $R_t$,
the structural UKF forms points
\[
  x_t^{(j)}=F_\theta(a_t^{(j)}),
  \qquad
  z_t^{(j)}=h_\theta(x_t^{(j)}),
\]
and then computes
\begin{align}
  \hat x_t
    &=\sum_j w_j^{(m)}x_t^{(j)},&
  \hat z_t
    &=\sum_j w_j^{(m)}z_t^{(j)},\\
  P_{xx,t}
    &=\sum_j w_j^{(c)}
      (x_t^{(j)}-\hat x_t)(x_t^{(j)}-\hat x_t)^\top,\\
  S_t
    &=\sum_j w_j^{(c)}
      (z_t^{(j)}-\hat z_t)(z_t^{(j)}-\hat z_t)^\top + R_t,\\
  P_{xz,t}
    &=\sum_j w_j^{(c)}
      (x_t^{(j)}-\hat x_t)(z_t^{(j)}-\hat z_t)^\top .
    \label{eq:bf-structural-ukf-moment-objects}
\end{align}
The Gaussian UKF update then uses
\begin{align}
  v_t &= y_t-\hat z_t,&
  K_t &= P_{xz,t}S_t^{-1},\\
  \hat x_{t|t}
    &= \hat x_t+K_tv_t,&
  P_{t|t}
    &= P_{xx,t}-K_tS_tK_t^\top,
    \label{eq:bf-structural-ukf-update-objects}
\end{align}
with log-likelihood contribution
\begin{equation}
  \ell_t
  =
  -\frac12
  \left[
    d_y\log(2\pi)+\log\det S_t+v_t^\top S_t^{-1}v_t
  \right].
  \label{eq:bf-structural-ukf-loglik-object}
\end{equation}
Thus a collapsed SVD route is not merely changing an implementation detail.
It replaces the point cloud
$\{F_\theta(a_t^{(j)})\}$ in
\eqref{eq:bf-structural-ukf-moment-objects} by
$\{\widetilde F_\theta(\widetilde a_t^{(j)})\}$.  The error enters through
\begin{align}
  \Delta\hat z_t
    &= \widetilde{\hat z}_t-\hat z_t,\\
  \Delta S_t
    &= \widetilde S_t-S_t,\\
  \Delta P_{xz,t}
    &= \widetilde P_{xz,t}-P_{xz,t},\\
  \Delta K_t
    &= \widetilde P_{xz,t}\widetilde S_t^{-1}
      -P_{xz,t}S_t^{-1}.
    \label{eq:bf-collapsed-ukf-object-errors}
\end{align}
Even when the formal sigma-point rule is accurate for its own input law,
\eqref{eq:bf-collapsed-ukf-object-errors} is controlled by whether the
collapsed cloud matches the structural cloud on the particular functions
$h_\theta$, $x\,h_\theta(x)^\top$, and
$h_\theta(x)h_\theta(x)^\top$ used in the UKF update.  Equality of predictive
laws is the robust sufficient condition; accidental equality of a few low-order
moments is not a model-contract proof.  This is the precise UKF sense in which
the law-level diagnostics above are not optional metadata.

The six diagnostics above correspond to the UKF recursion as follows.
\begin{enumerate}[label=(\roman*)]
  \item \textbf{Pushforward equality} is the condition that the two point clouds
    approximate the same expectations in
    \eqref{eq:bf-structural-ukf-moment-objects}.  If
    \eqref{eq:bf-collapsed-law-equivalence} fails, then
    $\widetilde{\hat x}_t$, $\widetilde{\hat z}_t$,
    $\widetilde S_t$, and $\widetilde P_{xz,t}$ are moments of a different
    predictive law.
  \item \textbf{Identity residuals} test whether every propagated point used in
    \eqref{eq:bf-structural-ukf-moment-objects} is an admissible structural
    state.  If
    $\Delta_{\mathrm{id}}^{\mathrm{col}}>0$, then the UKF averages include
    states that the DSGE transition assigns probability zero.
  \item \textbf{Point-count overhead} affects the number of evaluations of
    $F_\theta$ and $h_\theta$ required to build the same objects
    $\hat x_t,\hat z_t,S_t,P_{xz,t}$.  Extra zero-noise directions do not add
    information; they add factorization and map-evaluation work before the
    update in \eqref{eq:bf-structural-ukf-update-objects}.
  \item \textbf{Normal-support leakage} enters the UKF through $S_t$ and
    $P_{xz,t}$.  Artificial variance in a deterministic-completion direction
    changes the innovation covariance and the state-observation cross
    covariance, so it changes both the likelihood
    \eqref{eq:bf-structural-ukf-loglik-object} and the gain
    $K_t=P_{xz,t}S_t^{-1}$.
  \item \textbf{Spectral-gap gradient risk} is the derivative of the same UKF
    likelihood and update with respect to parameters.  If
    $\partial_\theta\widetilde C_t$ is unstable, then the gradients of
    $\widetilde S_t$, $\widetilde P_{xz,t}$, $\widetilde K_t$, and
    $\widetilde\ell_t$ inherit that instability.
  \item \textbf{Diagnostic separation} tells the reader whether
    \eqref{eq:bf-structural-ukf-update-objects} and
    \eqref{eq:bf-structural-ukf-loglik-object} were evaluated under the
    structural law or under a regularized law.  Without the identity and
    leakage diagnostics, a finite UKF likelihood cannot distinguish those two
    cases.
\end{enumerate}

A one-dimensional observation example shows the mechanism without requiring
filtering expertise.  Let
\[
  m_t=\varepsilon_t,\qquad k_t=m_t,\qquad
  y_t=k_t+e_t,\qquad
  \varepsilon_t\sim\calN(0,q),\quad e_t\sim\calN(0,r).
\]
The structural predictive state has
\[
  \Var(m_t)=q,\qquad
  \Var(k_t)=q,\qquad
  \Cov(m_t,k_t)=q,
  \qquad
  S_t^{\mathrm{str}}=q+r.
\]
The Kalman gain components implied by the UKF moment objects are therefore
\[
  K_t^{\mathrm{str}}
  =
  \frac{\Cov((m_t,k_t),y_t)}{\Var(y_t)}
  =
  \frac{1}{q+r}
  \begin{pmatrix}q\\ q\end{pmatrix}.
  \label{eq:bf-structural-gain-identity-example}
\]
If the collapsed SVD route adds a small independent numerical variance
$\eta^2$ to the deterministic-completion coordinate $k_t$, then
\[
  \widetilde k_t=m_t+\eta\xi_t,\qquad
  \xi_t\sim\calN(0,1),
\]
and the UKF objects become
\[
  \widetilde S_t=q+\eta^2+r,\qquad
  \widetilde K_t
  =
  \frac{1}{q+\eta^2+r}
  \begin{pmatrix}q\\ q+\eta^2\end{pmatrix}.
  \label{eq:bf-collapsed-gain-identity-example}
\]
The likelihood and update have changed:
\begin{align}
  \widetilde\ell_t-\ell_t
  &=
  -\frac12
  \left[
    \log\frac{q+\eta^2+r}{q+r}
    +
    y_t^2
    \left(\frac{1}{q+\eta^2+r}-\frac{1}{q+r}\right)
  \right],\\
  \widetilde K_{m,t}-K_{m,t}
  &=
  -\frac{q\eta^2}{(q+r)(q+\eta^2+r)},\\
  \widetilde K_{k,t}-K_{k,t}
  &=
  \frac{\eta^2 r}{(q+r)(q+\eta^2+r)} .
  \label{eq:bf-collapsed-ukf-likelihood-gain-change}
\end{align}
In this example, the six diagnostics above appear in concrete form:
\begin{enumerate}[label=(\roman*)]
  \item \textbf{Pushforward equality} is the difference between the structural
    law $k_t=m_t$ and the collapsed law
    $\widetilde k_t=m_t+\eta\xi_t$.  If $\eta>0$, the predictive law is no
    longer supported on $k=m$, and the UKF moments in
    \eqref{eq:bf-collapsed-gain-identity-example} are not the moments in
    \eqref{eq:bf-structural-gain-identity-example}.
  \item \textbf{Identity residuals} are measured by
    $c(m,k)=k-m$.  The structural points have $c=0$ pointwise, while the
    collapsed points have
    $c(\widetilde m_t,\widetilde k_t)=\eta\xi_t$ and therefore positive
    residual whenever the artificial direction is sampled.
  \item \textbf{Point-count overhead} is hidden in the construction of
    $\widetilde k_t$.  The structural rule needs points for the one genuine
    shock $\varepsilon_t$; the collapsed route adds a second numerical
    direction $\xi_t$ before computing the same scalar observation $y_t$.
  \item \textbf{Normal-support leakage} is the extra term $\eta^2$ in
    $\widetilde S_t=q+\eta^2+r$ in
    \eqref{eq:bf-collapsed-gain-identity-example}.  That term is exactly the
    artificial variance normal to the structural line $k=m$, and it changes the
    likelihood and gain in
    \eqref{eq:bf-collapsed-ukf-likelihood-gain-change}.
  \item \textbf{Spectral-gradient risk} enters if $\eta^2$ is produced by an
    SVD threshold, nugget branch, or spectral factor of a nearly singular
    covariance.  Then derivatives of
    $\widetilde S_t$, $\widetilde K_t$, and $\widetilde\ell_t$ inherit the
    derivative behavior of that spectral branch.
  \item \textbf{Diagnostic separation} is the need to report whether the
    likelihood is based on $S_t^{\mathrm{str}}=q+r$ or on
    $\widetilde S_t=q+\eta^2+r$.  Both are finite UKF likelihoods, but
    \eqref{eq:bf-collapsed-ukf-likelihood-gain-change} shows that they are not
    the same likelihood unless $\eta=0$.
\end{enumerate}
So the nugget changes both the likelihood weight assigned to the observation
and the way the observation is fed back into $m_t$ and $k_t$.  This is the UKF
problem in concrete terms: the filter can remain numerically stable and still
perform the prediction, likelihood evaluation, and update for a different
state-space model.

The practical conclusion is therefore conservative.  SVD is a good numerical
backend for singular or nearly singular covariance factors, and it can be used
inside a structurally correct sigma-point filter.  But using SVD on a collapsed
full-state covariance is usually a harder implementation contract: it has more
linear algebra to audit, more fragile derivatives for HMC, and more ways for a
regularization choice to become an unrecorded model change.  The structural
adapter is preferred because it states the economic law directly and leaves SVD
to do the narrower job of factoring the covariance that actually belongs to
that law.

For HMC promotion, one more distinction matters.  Matrix backpropagation
formulas for spectral decompositions contain denominators involving
singular-value or eigenvalue gaps, so repeated or nearly repeated spectral
values are derivative-risk regions \citep{ionescu2015matrixbackprop}.  That
risk is orthogonal to the structural-law issue.  A structurally correct SVD
sigma-point value path still needs spectral-gap telemetry, finite-gradient
stress tests, and compiled parity before it can be treated as a production HMC
target.

\section{Pruned DSGE and Adapter Implications}
\label{sec:bf-pruned-dsge-correct-pattern}

For pruned second-order DSGE approximations, the source DSGE project already
uses the standard first-order/second-order decomposition:
\begin{align}
  x_t^f &= h_x x_{t-1}^f + \eta\varepsilon_t,\\
  x_t^s &= h_x x_{t-1}^s
    + \frac{1}{2}H_{xx}(x_{t-1}^f\otimes x_{t-1}^f)
    + \frac{1}{2}h_{ss},\\
  y_t &= g_x(x_t^f+x_t^s)
    + \frac{1}{2}G_{xx}(x_t^f\otimes x_t^f)
    + \frac{1}{2}g_{ss} + e_t .
\end{align}
That decomposition is necessary, but it is not sufficient for large DSGE
filtering.  It separates perturbation order; it does not automatically
separate exogenous shock states from endogenous deterministic states inside
$x_t^f$.  A robust BayesFilter DSGE adapter should therefore carry both
bookkeeping axes,
\[
  \text{perturbation order: } (x^f,x^s),
  \qquad
  \text{structural timing: } (m,k).
\]
The adapter-facing implication is simple: the model must expose the exogenous
transition, the deterministic endogenous transition, the observation map, and a
declaration of which variables define the stochastic integration space.
For a sigma-point implementation, that adapter contract should compile to the
following concrete objects:
\begin{enumerate}[label=(\roman*)]
  \item dimensions and names for the filtered state, the declared innovation,
    and deterministic-completion coordinates;
  \item a function that maps each point
    $(x_{t-1}^{(j)},\varepsilon_t^{(j)})$ to
    $x_t^{(j)}=F_\theta(x_{t-1}^{(j)},\varepsilon_t^{(j)})$;
  \item a deterministic-identity residual, evaluated on every propagated point;
  \item an observation map evaluated only after structural completion;
  \item an approximation label if any legacy full-state or regularized
    covariance path is used instead.
\end{enumerate}
The pruned perturbation decomposition tells the adapter what equations exist;
the structural metadata tells the filter which variables receive quadrature.

\section{Source-Project Lesson}
\label{sec:bf-current-source-project-failure-mode}

The recent source-project audit found a concrete instance of the structural
failure discussed in this chapter.  The default DSGE SVD sigma-point adapter
uses a collapsed transition of the form
\[
  Q_w = \eta\eta^\top + \epsilon I,\qquad
  x_{t|t-1}^f = h_x x_{t-1|t-1}^f,
\]
and propagates the covariance as
\[
  P_{t|t-1} = h_x P_{t-1|t-1} h_x^\top + Q_w.
\]
It tracks the pruned $x^s$ correction separately, but it does not expose a
partition between exogenous and endogenous coordinates inside $x^f$.  For the
smallest NK example, where the states are only the exogenous demand and
monetary-policy shocks, this is not a serious modeling error.  For Rotemberg NK,
SGU, EZ, and NAWM-scale DSGE models, however, the state vector contains
endogenous/predetermined components, so treating every coordinate as an
exchangeable noisy Gaussian coordinate becomes a structural filtering bug rather
than a mere numerical tuning issue.

This does not mean the generic SVD sigma-point machinery is unusable.  The
source audit found a more precise split: a generic sigma-point path that accepts
user-supplied transition points or a structural transition map can be reused as
an approximate backend, while the default mixed-DSGE adapter must be gated
unless it exposes the exogenous/endogenous partition and deterministic
completion.  In other words, reuse the factorization and diagnostics; do not
reuse an adapter that silently changes the stochastic integration space.

\section{Validation Gates and Final Policy Rule}
\label{sec:bf-structural-deterministic-required-tests}

No future BayesFilter report should claim that a nonlinear DSGE filter is
value-valid, gradient-valid, sampler-usable, converged, or production-ready in
the sense of Chapter~\ref{ch:bf-production-checklist} unless the following
structural tests are passed or explicitly waived with written justification:
\begin{enumerate}
  \item \textbf{Metadata test:} the model exposes exogenous, endogenous, and
    deterministic state blocks, and their dimensions match the shock-impact
    matrix and transition maps.
  \item \textbf{Constraint-support test:} propagated sigma points or particles
    satisfy deterministic identities to numerical tolerance.
  \item \textbf{Linear recovery test:} when the structural transition is
    linear, the structural nonlinear filter recovers the exact Kalman
    likelihood up to the declared approximation tolerance.
  \item \textbf{Degenerate-transition test:} zero-innovation endogenous rows
    remain deterministic except for explicitly declared numerical
    regularization.
  \item \textbf{SVD/factorization test:} if a square-root or SVD sigma-point
    backend is used, diagnostics state whether the factor belongs to the
    structural integration variable or to a legacy full-state approximation;
    the latter must carry an approximation label.
  \item \textbf{Toy DSGE test:} a controlled model with
    \[
      m_t = \rho m_{t-1}+\sigma\varepsilon_t,\qquad
      k_t=f(k_{t-1},m_{t-1},m_t),
    \]
    compares structural propagation against a dense quadrature or
    high-particle reference.
  \item \textbf{Case-study test:} Rotemberg NK, SGU, and any EZ/NAWM target
    must document which state coordinates are shock-driven, which are
    deterministic conditional on the shock path, and how the filter honors
    that timing.
  \item \textbf{Derivative-promotion test:} any HMC claim through SVD or
    eigenspectral factors must pass spectral-gap, finite-gradient, and
    eager/compiled parity checks.
\end{enumerate}

\subsection{Common Misunderstandings}
\label{sec:bf-structural-common-misunderstandings}

The most common reader errors in this area are:
\begin{enumerate}[label=(\roman*)]
  \item a deterministic endogenous state need not have zero one-step predictive
    variance; it can inherit randomness through the declared stochastic block;
  \item a singular or rank-deficient transition covariance is not, by itself, a
    modeling error in the exact linear-Gaussian case;
  \item numerical regularization is not the same thing as declaring a new
    stochastic state coordinate;
  \item a full-state nonlinear approximation is not automatically forbidden, but
    it must be labeled as an approximation if it enlarges the stochastic space
    beyond the structural transition contract.
\end{enumerate}

\begin{center}
\fbox{\parbox{0.88\linewidth}{
For nonlinear DSGE filtering, the stochastic integration variables are not
automatically the entries of the state vector.  They are the variables declared
stochastic by the structural transition.  Deterministic endogenous states must
be computed, not noised.
}}
\end{center}

If an implementation violates this rule, it must be labeled as an artificial
approximation and justified against a reference.  Otherwise, the filter is
solving a different approximate state-space model relative to the declared
structural law.  Failing one of the validation gates above does not
automatically make a backend useless, but it does mean the strongest claim
labels from Chapter~\ref{ch:bf-production-checklist} are not yet justified.
